%& -shell-escape
\documentclass[12pt,%                     % corpo del font principale
               a4paper,twoside,		      % carta A4, fronte
               titlepage,%                % frontespzio
               headinclude,,footinclude,% % testatina e pié di pagina
               numbers=noenddot,%         % niente punto dopo il numero delle sezioni
               cleardoublepage=empty,%    % pagine vuote senza testatina e pié di pagina
               abstracton,%
               appendixprefix=true,%
               %draft                     % useful for debug
               ]{scrreprt}                % classe report di KOMA-Script;

%%% font / input %%%%%%%%%%%%

\usepackage[T1]{fontenc}                  % la codifica dei font
\usepackage[utf8]{inputenc}               % imposta la codifica di input;

%%% language %%%%%%%%%%%%%%%%

\usepackage[italian,english]{babel}       % per scrivere in italiano e in inglese;

%%% document style %%%%%%%%%%

\usepackage[eulerchapternumbers,%         % numeri dei capitoli in Euler
            subfig,%                      % compatibilità con subfig
            eulermath,%                   % AMS Euler come font per la matematica
            %pdfspacing%                  % migliora il riempimento di riga con PDFLaTeX
            ]{classicthesis}              % lo stile ClassicThesis
\usepackage{arsclassica}                  % modifica alcuni aspetti di ClassicThesis
\usepackage{setspace}                     % setting space among lines
\onehalfspacing
\KOMAoptions{cleardoublepage=plain}       % the cleardoublepage is empty

%%% bibliography %%%%%%%%%%%%

\usepackage[autostyle]{csquotes}          % useful with biblatex
\usepackage[backend=biber,                % bibliography
            bibstyle=alphabetic,          % .
			citestyle=alphabetic,         % .
			hyperref,                     % .
			useprefix,                    % .
			backref]{biblatex}            % .
\addbibresource{refs/bibliografia.bib}    % loads bib file

\DeclareBibliographyCategory{web}
\addtocategory{web}{wiki:hull}
\defbibheading{web}{\subsection*{Online resources}}

%%% math and science %%%%%%%%%

\usepackage{amsmath,amssymb,amsthm}       % indispensabili per la matematica
\usepackage{xfrac}                        % slanted bar for fractions
\usepackage{siunitx}	                  % SI units
\usepackage{eurosym}                      % euro symbol

%%% layout %%%%%%%%%%%%%%%%%%

\usepackage[{a4paper,                     % margins
             left=3cm,                    % .
			 right=2.5cm,                 % .
             top=3cm,                     % .
			 bottom=3.5cm,                % .
             heightrounded,               % .
			 bindingoffset=1.5cm}]{geometry}
\usepackage{varioref}                     % riferimenti completi della pagina
\usepackage[plain]{fancyref}
%\usepackage{indentfirst}                 % rientra il primo paragrafo di ogni sezione
\usepackage{minitoc}
\usepackage{wrapfig}

%%% tables %%%%%%%%%%%%%%%%%%%

\usepackage{booktabs}			          % migliora la qualità delle tabelle
\usepackage{tabularx}			          % colonne a spaziatura fissa delle tabelle
\newcommand{\otoprule}                    % better top rule horizontal line
    {\midrule[\heavyrulewidth]}           % .

%%% images and graphics %%%%%%

\graphicspath{{img/}}                     % defines images path
\usepackage{graphicx}                     % basic graphics support
\usepackage{pdfpages}                     % include pdf
\usepackage{subfig}                       % sottofigure, sottotabelle
\usepackage{sidecap}                      % enables captions on the side
\usepackage[tableposition=top]{caption}   % define captions style
\usepackage{float}
\usepackage{parskip}

%% tikz
\usepackage{dot2texi}
\usepackage{tikz}
\usetikzlibrary{shapes,arrows}
\usepackage{gnuplot-lua-tikz}
\usepackage{rotating}
\usepackage[dvipsnames]{xcolor}

\newcommand\AlCentroPagina[1]{            % generates big b/w logo
    \AddToShipoutPicture*{\AtPageCenter{  % .
        \makebox(0,0){\includegraphics    % .
            [width=0.9\paperwidth]{#1}}   % .
        }                                 % .
    }                                     % .
}                                         % .

%%% text utils %%%%%%%%%%%%%%%

\usepackage{hyphenat}
\hypersetup{colorlinks=false}		      % commentare se il file non è x la stampa
\usepackage{acronym}			          % pacchetto per gli acronimi

% generates summary (before starting chapter's first section)
\newenvironment{chaptersum}
    {\sffamily\textbf{Sommario}\\\rightskip1in}
    {}


\begin{document}
    \includepdf[pages={1}]{pages/tesi-front}
    \cleardoublepage
    \includepdf[pages={1}]{pages/tesi-post-front}
    \cleardoublepage
    %******************************************
    % Materiale iniziale
    %******************************************
    \pagenumbering{Roman}
    \thispagestyle{empty}
\begin{flushright}
    \null\vspace{\stretch {1}}
        To my family
    \vspace{\stretch{2}}\null
\end{flushright}
\cleardoublepage
    \pagestyle{scrheadings} 
    \pdfbookmark{\contentsname}{tableofcontents}
\setcounter{tocdepth}{2}
\dominitoc\tableofcontents

\pdfbookmark{\listfigurename}{lof}
\listoffigures

\pdfbookmark{\listtablename}{lot}
\listoftables
\cleardoublepage
    %\chapter*{Forewords}

    some nice fore text
\cleardoublepage
    %******************************************
    % Materiale principale
    %******************************************
    \pagenumbering{arabic}
    \chapter{Introduction}

    \begin{chaptersum}
        Testo sommario contenuti in italiano
    \end{chaptersum}

    The history of non-invasive geophysical research and methods has evolved rapidly over the years, widening the horizons of research in all the fields in which it has been applied. In this context, archaeology is one of the disciplines which has taken advantage of this and the number of published experiences on the topic is growing every year.\\
    In the last 20 years, both geophysics and archaeology have steadily expanded their use of Geographic Information Systems (GIS) to organize, clean, analyze and produce results from raw data, and 35 years have past since the first application of GIS in archaeology \cite[pp.~2--3]{kvamme1995}. From this point of view, GIS placed itself in the intersection of both fields and has developed itself as a separate discipline over the years, moving from desktop software to the internet with \emph{webGIS}. The ever growing amount of case studies focused on GIS, geophysics and archaeology is making clear that in the future these disciplines will evolve together towards an interdisciplinary approach, and probably each of these will take advantage from the advancements in the others.\\
    This works presents a case study of a statistical approach on geophysical data recorded in the archaeological field, using some well known methods and introducing some innovations with respect to data analysis, webGIS and interfaces to present results. 
    \section{Archaeology and geophysics}
        The relationship between geophysics and \emph{landscape archaeology} in particular is well documented and presents some practical advantages over any other approach, mainly:
        \begin{description}
            \item[cost reduction] geophysical surveys can drive the excavation and avoid losing time in unfruitful diggings, helping to define the best approach to manage the site
            \item[flexible approaches] geophysical methods are suitable for large scale analysis as well as small scale surveys
            \item[integration] since part of the archaeological data consists of spatial data, geophysical information --- most of the times --- could be easily integrated with current data, defining another ``layer of knowledge''.
        \end{description}

    \section{Landscape archaeology and Geographic Information Systems}
        The natural tendency of modern archaeology to move from small scale excavations to a larger view and comprehension of the relationship between ancient communities and past natural environment --- the so called \emph{landscape archaeology} --- has brought new methods in archaeology, and most of these relies on Geographic Information Systems. In this field, this relationship has gained strength more than with any other archaeological approach, as most of the landscape archaeology methods beyond excavation are \emph{spatial methods} \cite[preface]{space-archaeology}. The commonest uses of GIS in this field are: data classification, data plotting, environment modeling; the first two have been used in this work.

        By definition \cite{american-dict}, the scientific method consists of
        \begin{quote}
            the observation, identification, description, experimental investigation, and theoretical explanation of phenomena. Such activities restricted to a class of natural phenomena.
        \end{quote}
        Landscape archaeology deals with spatially described phenomena, and in modern approaches, GIS are the best tool to manage, identify and do experimental investigation on phenomena distributed on everything more complex than a Cartesian surface (like a geoid).

    \section{Geophysics analysis workflow}
        Despite of some peculiarities depending on the approach, the workflow of a geophysical analysis in the archaeological context has a fixed structure consisting of a sequence of data collecting techniques, cleaning of the collected data and GIS methods to clean, calculate and plot results. These steps, with particular attention to topics related to the current work, are briefly described below.

        \subsection{Aerial photography}
            The term \emph{aerial} refers to any image of the ground taken from an elevated position \cite{wiki:aerial}; in geophysics it is used to describe the images taken from aircrafts, helicopters or balloon with the specific purpose to define the characteristics of the ground before starting a deeper investigation with more precise methods. It is useful to have a general idea of where heritage is on the ground and its relations with the natural and anthropic environment.\\
            With reference to the nature of the work that will be presented in the following chapters, the aerial (oblique) survey gives best results in a plain farmed ground, providing effective documentation on cropmarks, classical and medieval structures (as well as their place in rural or urban landscape) \cite[pp.~11--12]{arch-site-detection}. Before satellite imaging, it has been one of the most valuable survey techniques and its use rapidly matured during the two World Wars.
            
        \subsection{Magnetometry}
            Magnetic techniques are proved as the most effective ones, when the necessity rises to define the position and size of buried structures. In its basic principles, it measures the anomalies on the measured Earth's magnetic field on the ground, caused by the presence of buried archaeological remains. Different configurations of the magnetometer's sensors are available (e.g. \emph{duo--quadro} sensors) and usually it is operated moving sensors on parallel equispaced lines (\num{0.1} to \SI{0.5}{\meter}), enabling the survey of large areas in a very small amount of time. The obtained data must be processed by specialized software (cleaning them by defects like \emph{spikes}, \emph{stripes} and \emph{zig-zag}) to obtain meaningful representations of the measures \cite[p.~47]{remote-ciminale}; these representations give a detailed visual insight and 3D appearance of the underlying structures.

        \subsection{GIS and vector data creation}
            Analyzing data from different sources and with different contents like photos and magnetic prospecting would be impossible without the help of a spatial analysis software, which is in fact merging them in a single entity (the GIS project) using the only common attribute: the geographical coordinates. The data (after being georeferred and projected with a suitable geographical projection) are stored as \emph{raster} (actually, images with a latitude--longitude tuple associated with every pixel) in a plain file or database.\\

            \subsubsection{Deriving vectors from observation}
                Having as base layers the information-rich rasters, the operator can draw GIS geometries (\emph{points}, \emph{lines} or \emph{polygons}) to represent the features visible in the data. Notably, this is an interpretation itself: the operator is creating new knowledge starting from available data from different sources, assigning attributes to each feature to identify it (\emph{id}) or to classify it (\emph{road}, \emph{building}, \emph{ditch}). Moreover, geographical attributes are automatically stored in the newly created geometries, and all the GIS software are able to export them in common exchange formats (\emph{ESRI shapefile}, \emph{geoJSON}, etc.).

            \subsubsection{Deriving vectors from analysis}
                Most GIS suites integrate sets of functions calling algorithms to derive information from available data. E.g., a perimeter can be derived from a subset of all the features, or a buffer geometry can be created from a river basin's perimeter. All these data are geometries and could be saved in a separate layer, enriching the data collection (literature offers a great variety of these operations in a geophysical context \cite[p.~325]{remote-ciminale}). Depending on the GIS software used, functions may be implementations of well known --- published --- algorithms or written by the software's producer and distributed as \emph{proprietary} algorhitms; the latter case is the less advisable for the user, since he will know the results of the calculation without knowing the how the calculation itself has been operated \cite[p.~69]{fronza-informatica}.

        \subsection{GIS data management\label{sec:gis-data-management}}
            The workflow defined in this chapter has some points of weakness; some of them are related to the operator when deriving vectors from observation, others are limitations of the software suite itself:
            \begin{description}
                \item[errors related to magnification] when the GIS user derives vector data (usually polygons) from magnetic prospections, the zoom level set while editing the layer's features is fundamental for an accurate drawing; if the base layers are set at a zoom level lower than the adequate one, a rising of the difference between derived geometries and original raster layers will take place, heavily decreasing the accuracy and precision of the new data.
                \item[database management] collecting all the abovementioned data types for hundreds of sites will make the database grow heavily, raising the necessity to have an efficient data storage architecture. In the geospatial field, this usually means to stop using ESRI Shapefiles in favour of spatially-enabled database (Spatialite, PostGIS, etc.). The approach to these technologies is not simple most of the times, and lowering the barriers to access geospatial functions in database becomes a priority.
                %\item[results exportation] 
                \item[desktop software] most of the analysis process is usually done using functions integrated in a desktop GIS software, making very easy to move data across physical storage devices, but decreasing consistently the data safeness (since data needs to be backed up regularly and the machine itself can be damaged by other people in the laboratory). Moreover, giving external collaborators access to desktop data is very difficult, as well as to manage --- or log --- the access to data from a single person. These disadvantages are well documented in literature (see \cite[p.~19]{fronza-informatica}).
            \end{description}
\cleardoublepage
    \chapter{Case study: geostatistics on neolithic settlements in Tavoliere}

    \begin{chaptersum}
        Sommario in italiano
    \end{chaptersum}

    The geostatistical system described in this chapter has been structured using the processed spatial data collected during almost ten years of surveying of the settlements in the Tavoliere plain. The existence of buried settlements ranging from neolithic period to middle ages has been proved by field and aerial surveys and by historical records. Given the variety and widespread distribution of the settlements across about \SI{150}{\hectare}, aerial and geophysical methods assured the best results, and heavily contributed to collect the data useful to an accurate landscape archaeology study \cite[pp.~45--48]{remote-ciminale}.

    \section{General features of the settlements}
        As extensively reported in literature \cite{intro-tavoliere}, the 256 sites form one of the densest concentrations of prehistoric settlements in Europe, lying within a plain approximately \SI{50}{\kilo\meter} by \SI{80}{\kilo\meter} at its broadest and longest points. The natural boundaries of this area are marked from rivers Fortore and Ofanto --- north and south --- and the Gulf of Manfredonia and Appennine foothills --- east and west.\\
        The settlements are large or small villages, normally surrounded by one or more ditches, with the interior generally filled with a number of internal compounds. The majority of the sites occupied level ground. Refer to published materials for further details.

        \subsection{Published data and the work of J. D. B. Jones}

            \subsubsection{Spatial analysis}

        \subsection{Modern spatial investigations}
            \begin{table}
    \centering
    \begin{tabular}{cccc}
        \toprule
        Layer type & Description          & Features & Scale \\
        \otoprule
        raster     & magnetic prospection & ---      & TODO  \\
        raster     & aerial photography   & ---      & TODO  \\
        vector     & ditches / compounds / areas & 1300   & ---   \\
        \bottomrule
    \end{tabular}
    \caption[List of layers in Tavoliere neolithic GIS project]{Details of the layers in the Tavoliere GIS project.
    \label{tab:layers}
    }
\end{table}
% tif con griglia magnetogramma
% magnetogramma tif
% foto ortorettificata
% layer vect punti GPS estremi magnetogramma

% viene prima inserito magnetogramma con griglia in unico file raster, poi derivato magnetogramma senza griglia
% georeferenziazione viene operata su magnetogramma con griglia

% WMS SIT Puglia
% WMS Geoportale Nazionale
% cartografia IGM 25000
% carta geologica

            After the post-processing, all the data have been placed in a GIS environment, currently containing XXX different raster layers and a single vector layer (\fref{tab:layers}). These vector layers refer to the latest study in chronological order, carried out by dott.ssa Angela Laterza in 2013,
            % TODO: add bib reference to Angelica's thesis
            who completely digitized the structures contained in 23 neolithic settlements. The whole digitalization process took as long as 320 hours and consisted in manually tracing the borders of the visible structures (\emph{ditches} and \emph{compounds}) and, as a separate geometry, the respective area for each of them; the operator has been employed for about two months, and has produced about 1300 geometries; \fref{fig:scheme-derive} shows the whole process.

            \begin{figure}[htb]
                \resizebox{1.1\textwidth}{!}{%
                    \input{tab/tab-geo-workflow}
                }
                \caption[Data deriving workflow for the Tavoliere project]{Data derivation workflow for the case study. Raster layers (\textsf{rs}) are the source of all derived vector data (\textsf{vt}). Bold lines represents a manual process, while dashed ones an automated process. Perimeter as numeric value have been calculated and saved in derived geometries.}
                \label{fig:scheme-derive}
            \end{figure}

    \section{The GIS approach to statistics}
        The first step in statistical analysis at any level is to build a suitable set of data. In this case study, vector shapes of ditches and compounds can be considered as \emph{primary} data, while area geometries as \emph{derived} data.
        The statistical and spatial approach to the study of neolithic --- or, in general, prehistoric --- settlements data is well documented in literature \cite{arch-location-model}, and the strict relation between the morphology of the natural environment and the ancient societies enforces the use of GIS systems when deriving data.\\
        Moving from the standard geophysical survey workflow to a spatial analysis context during the study of the settlements has generated some necessities, which need to be added to the critical points already defined in \fref{sec:gis-data-management}:

        \begin{description}
            \item[automate derived data creation] the set of derived data may be generated automatically, since it consists of geometries describing the same characteristic on every shape (e.g. the internal area of a ditch, for all ditches);
            \item[automate calculations] most of the processes operated on data take a feature of the element (e.g. the perimeter), save it in a database field, retrieves all of them and make calculations; this is a repetitive task which could be done more efficiently;
            \item[platform independent results] the results must be reproducible by any researcher on any platform, and the dynamics of calculations must be public; the structure of the project should enable other researchers to contribute and improve the algorithms.
        \end{description}

        Given that, some objectives have been defined to improve the current workflow and speed up the process, trying to gain in the same time more precision, accuracy and reproducibility; the following subsections will explain how every objective has been achieved in detail.

        \subsection{Development of the open source webGIS interface}

        \subsection{Bulk distinguishing ditches and compounds}
            In the current approach, ditches and compounds geometries are distinguished by attributes (the \textsf{Ditch\_comp} field in the Shapefile is respectively set to $1$ or $2$), manually added by the operator to each structure in the Shapefile.

            \begin{wrapfigure}{l}{0.5\textwidth}
                \centering
                \begin{tikzpicture}[x=1mm,y=1mm,scale=0.005]
                    \input{tab/dot-flow-map}
                \end{tikzpicture}
                \caption[Flow chart: the logic of bulk distinguishing ditches and compounds]{If any of the ditches and compounds are sharing the same color (class), class total number $k$ must be changed. At the end, geometry type is saved as an attribute in the shapefile.}
                \label{fig:flow-map}
            \end{wrapfigure}

            This operation has been semi-automated using the Jenks Natural Breaks classification method \cite{jenks1977}. Virtually any classification could have been chosen (quantile, standard deviation), but Jenks Natural Breaks has gained popularity in the last 20 years as a tool for coloring map objects based on objects properties (\emph{choropleth maps}), and it is a well known and tested tool\footnote{Some criticism has recently raised around the use of Jenks Natural Breaks for the classification of \emph{all} kinds of data; alternative solutions have been found for some distributions --- as heavy-tailed ones \cite{jenks-tail} --- but this is not the case, and Jenks has been experimentally proved as the most appropriate choice.}. It aims to present a series of break values that best represent the actual breaks observed in the data as opposed to some arbitrary classificatory scheme (i.e. equal interval); in this way the actual clustering of data values is preserved.

            In the case study, geometries have been classified by perimeter values (calculated and saved automatically from ditches and compounds shapes). The different coloring of the geometries on the map based on the newly created classes (\fref{fig:jenks-color}) helps the user to distinguish ditches from compounds at a glance.

            The process is semi-automated since the user have to manually select the ditches and, if necessary, change the class numbers to have a good fit. The flowchart in \fref{fig:flow-map} reports the logic of the process. At the end, the geometry types are saved as text attribute (literally \textsf{ditch} or \textsf{compound}) in a table containing the results of the processing (one row for each geometry); \fref{tab:jnb-results} shows sample data from the \emph{Anglisano} settlement.

            \begin{figure}[H]
                \centering
                \rotatebox{180}{
                    \begin{tikzpicture}[x=1mm,y=1mm,scale=0.2]
                        \definecolor{cff0000}{RGB}{207,240,0}
\definecolor{c2b83ba}{RGB}{194,184,59}
\definecolor{cabdda4}{RGB}{202,189,118}
\definecolor{cffffbf}{RGB}{207,255,251}
\definecolor{cfdae61}{RGB}{207,218,230}
\definecolor{cd7191c}{RGB}{205,113,145}

\begin{scope}[draw=black,fill=cd7191c,line join=round,line width=0.208pt]
  \path[draw,fill] (545.0000,190.0000) -- (544.0000,191.0000) --
    (543.0000,191.0000) -- (543.0000,190.0000) -- (543.0000,189.0000) --
    (544.0000,188.0000) -- (545.0000,187.0000) -- (546.0000,187.0000) --
    (548.0000,181.0000) -- (548.0000,180.0000) -- (548.0000,179.0000) --
    (548.0000,178.0000) -- (547.0000,177.0000) -- (547.0000,176.0000) --
    (546.0000,176.0000) -- (547.0000,176.0000) -- (547.0000,175.0000) --
    (548.0000,175.0000) -- (549.0000,176.0000) -- (550.0000,178.0000) --
    (550.0000,180.0000) -- (550.0000,182.0000) -- (549.0000,184.0000) --
    (548.0000,185.0000) -- (545.0000,190.0000);
  \path[draw,fill] (538.0000,191.0000) -- (537.0000,191.0000) --
    (536.0000,191.0000) -- (535.0000,191.0000) -- (534.0000,191.0000) --
    (533.0000,190.0000) -- (531.0000,190.0000) -- (531.0000,189.0000) --
    (529.0000,185.0000) -- (528.0000,182.0000) -- (528.0000,181.0000) --
    (528.0000,179.0000) -- (529.0000,178.0000) -- (530.0000,176.0000) --
    (531.0000,175.0000) -- (532.0000,173.0000) -- (532.0000,172.0000) --
    (533.0000,172.0000) -- (533.0000,171.0000) -- (534.0000,171.0000) --
    (535.0000,171.0000) -- (536.0000,171.0000) -- (535.0000,172.0000) --
    (533.0000,175.0000) -- (532.0000,175.0000) -- (531.0000,177.0000) --
    (530.0000,178.0000) -- (530.0000,180.0000) -- (530.0000,182.0000) --
    (530.0000,183.0000) -- (530.0000,184.0000) -- (531.0000,185.0000) --
    (532.0000,187.0000) -- (533.0000,188.0000) -- (534.0000,189.0000) --
    (535.0000,190.0000) -- (538.0000,191.0000);
  \path[draw,fill] (339.0000,356.0000) -- (338.0000,356.0000) --
    (338.0000,355.0000) -- (338.0000,354.0000) -- (339.0000,353.0000) --
    (339.0000,352.0000) -- (340.0000,351.0000) -- (341.0000,351.0000) --
    (342.0000,350.0000) -- (343.0000,350.0000) -- (344.0000,350.0000) --
    (345.0000,350.0000) -- (346.0000,350.0000) -- (347.0000,351.0000) --
    (349.0000,351.0000) -- (351.0000,352.0000) -- (352.0000,353.0000) --
    (353.0000,354.0000) -- (354.0000,355.0000) -- (355.0000,357.0000) --
    (355.0000,359.0000) -- (355.0000,360.0000) -- (355.0000,362.0000) --
    (355.0000,364.0000) -- (354.0000,365.0000) -- (354.0000,366.0000) --
    (353.0000,366.0000) -- (353.0000,367.0000) -- (352.0000,368.0000) --
    (352.0000,369.0000) -- (351.0000,369.0000) -- (350.0000,370.0000) --
    (349.0000,370.0000) -- (348.0000,370.0000) -- (347.0000,370.0000) --
    (347.0000,369.0000) -- (346.0000,369.0000) -- (345.0000,369.0000) --
    (344.0000,369.0000) -- (343.0000,368.0000) -- (342.0000,368.0000) --
    (341.0000,367.0000) -- (340.0000,365.0000) -- (339.0000,364.0000) --
    (343.0000,364.0000) -- (347.0000,364.0000) -- (347.0000,363.0000) --
    (348.0000,363.0000) -- (349.0000,362.0000) -- (349.0000,361.0000) --
    (349.0000,360.0000) -- (349.0000,359.0000) -- (349.0000,358.0000) --
    (349.0000,357.0000) -- (348.0000,357.0000) -- (348.0000,356.0000) --
    (347.0000,356.0000) -- (346.0000,355.0000) -- (345.0000,355.0000) --
    (344.0000,354.0000) -- (343.0000,354.0000) -- (342.0000,353.0000) --
    (341.0000,354.0000) -- (340.0000,354.0000) -- (340.0000,355.0000) --
    (339.0000,355.0000) -- (339.0000,356.0000);
  \path[draw,fill] (330.0000,344.0000) -- (342.0000,343.0000) --
    (343.0000,342.0000) -- (343.0000,341.0000) -- (344.0000,340.0000) --
    (344.0000,339.0000) -- (344.0000,337.0000) -- (344.0000,336.0000) --
    (344.0000,335.0000) -- (344.0000,334.0000) -- (343.0000,333.0000) --
    (342.0000,332.0000) -- (341.0000,331.0000) -- (340.0000,331.0000) --
    (334.0000,331.0000) -- (336.0000,326.0000) -- (341.0000,328.0000) --
    (347.0000,330.0000) -- (348.0000,331.0000) -- (349.0000,333.0000) --
    (349.0000,335.0000) -- (349.0000,336.0000) -- (349.0000,338.0000) --
    (349.0000,340.0000) -- (348.0000,342.0000) -- (347.0000,343.0000) --
    (345.0000,345.0000) -- (333.0000,347.0000) -- (330.0000,344.0000);
  \path[draw,fill] (424.0000,299.0000) -- (421.0000,292.0000) --
    (425.0000,289.0000) -- (429.0000,286.0000) -- (431.0000,286.0000) --
    (432.0000,285.0000) -- (437.0000,286.0000) -- (441.0000,287.0000) --
    (442.0000,288.0000) -- (445.0000,289.0000) -- (446.0000,290.0000) --
    (447.0000,291.0000) -- (448.0000,292.0000) -- (449.0000,294.0000) --
    (449.0000,295.0000) -- (449.0000,297.0000) -- (449.0000,298.0000) --
    (450.0000,299.0000) -- (450.0000,300.0000) -- (450.0000,301.0000) --
    (449.0000,302.0000) -- (448.0000,302.0000) -- (447.0000,302.0000) --
    (446.0000,302.0000) -- (446.0000,301.0000) -- (446.0000,300.0000) --
    (447.0000,295.0000) -- (447.0000,293.0000) -- (444.0000,291.0000) --
    (438.0000,288.0000) -- (436.0000,288.0000) -- (434.0000,288.0000) --
    (432.0000,288.0000) -- (430.0000,289.0000) -- (429.0000,290.0000) --
    (427.0000,291.0000) -- (427.0000,292.0000) -- (425.0000,295.0000) --
    (424.0000,298.0000) -- (424.0000,299.0000);
  \path[draw,fill] (380.0000,263.0000) -- (380.0000,261.0000) --
    (381.0000,260.0000) -- (383.0000,259.0000) -- (385.0000,259.0000) --
    (387.0000,259.0000) -- (389.0000,259.0000) -- (391.0000,259.0000) --
    (392.0000,260.0000) -- (393.0000,261.0000) -- (394.0000,262.0000) --
    (395.0000,262.0000) -- (396.0000,266.0000) -- (396.0000,271.0000) --
    (395.0000,276.0000) -- (395.0000,279.0000) -- (394.0000,280.0000) --
    (393.0000,282.0000) -- (391.0000,284.0000) -- (389.0000,285.0000) --
    (387.0000,286.0000) -- (384.0000,287.0000) -- (384.0000,288.0000) --
    (383.0000,288.0000) -- (383.0000,289.0000) -- (382.0000,289.0000) --
    (381.0000,289.0000) -- (380.0000,289.0000) -- (380.0000,288.0000) --
    (379.0000,288.0000) -- (379.0000,287.0000) -- (382.0000,285.0000) --
    (387.0000,283.0000) -- (392.0000,279.0000) -- (392.0000,278.0000) --
    (393.0000,274.0000) -- (394.0000,270.0000) -- (393.0000,266.0000) --
    (392.0000,262.0000) -- (389.0000,261.0000) -- (385.0000,261.0000) --
    (382.0000,262.0000) -- (380.0000,263.0000);
  \path[draw,fill] (334.0000,431.0000) -- (333.0000,431.0000) --
    (333.0000,430.0000) -- (333.0000,429.0000) -- (334.0000,429.0000) --
    (334.0000,428.0000) -- (335.0000,428.0000) -- (336.0000,428.0000) --
    (336.0000,429.0000) -- (337.0000,429.0000) -- (338.0000,429.0000) --
    (339.0000,429.0000) -- (340.0000,429.0000) -- (341.0000,430.0000) --
    (342.0000,431.0000) -- (342.0000,432.0000) -- (343.0000,433.0000) --
    (342.0000,434.0000) -- (343.0000,434.0000) -- (343.0000,437.0000) --
    (342.0000,440.0000) -- (341.0000,443.0000) -- (339.0000,444.0000) --
    (337.0000,444.0000) -- (336.0000,444.0000) -- (334.0000,443.0000) --
    (333.0000,443.0000) -- (333.0000,442.0000) -- (333.0000,441.0000) --
    (333.0000,440.0000) -- (334.0000,439.0000) -- (334.0000,438.0000) --
    (335.0000,438.0000) -- (336.0000,439.0000) -- (337.0000,439.0000) --
    (338.0000,439.0000) -- (339.0000,438.0000) -- (340.0000,438.0000) --
    (340.0000,437.0000) -- (341.0000,437.0000) -- (341.0000,436.0000) --
    (341.0000,435.0000) -- (340.0000,434.0000) -- (340.0000,433.0000) --
    (339.0000,432.0000) -- (338.0000,432.0000) -- (334.0000,431.0000);
  \path[draw,fill] (438.0000,204.0000) -- (435.0000,202.0000) --
    (438.0000,198.0000) -- (442.0000,196.0000) -- (445.0000,196.0000) --
    (447.0000,195.0000) -- (448.0000,195.0000) -- (451.0000,199.0000) --
    (453.0000,202.0000) -- (454.0000,204.0000) -- (454.0000,205.0000) --
    (454.0000,206.0000) -- (454.0000,207.0000) -- (454.0000,208.0000) --
    (453.0000,209.0000) -- (451.0000,209.0000) -- (449.0000,208.0000) --
    (450.0000,206.0000) -- (451.0000,203.0000) -- (450.0000,202.0000) --
    (450.0000,201.0000) -- (449.0000,200.0000) -- (448.0000,199.0000) --
    (446.0000,199.0000) -- (445.0000,198.0000) -- (438.0000,204.0000);
  \path[draw,fill] (486.0000,193.0000) -- (486.0000,192.0000) --
    (487.0000,192.0000) -- (490.0000,195.0000) -- (496.0000,200.0000) --
    (501.0000,207.0000) -- (502.0000,210.0000) -- (502.0000,211.0000) --
    (503.0000,212.0000) -- (503.0000,214.0000) -- (503.0000,215.0000) --
    (502.0000,217.0000) -- (502.0000,218.0000) -- (501.0000,219.0000) --
    (500.0000,220.0000) -- (499.0000,220.0000) -- (499.0000,221.0000) --
    (498.0000,221.0000) -- (497.0000,221.0000) -- (497.0000,220.0000) --
    (498.0000,220.0000) -- (499.0000,219.0000) -- (500.0000,217.0000) --
    (500.0000,215.0000) -- (501.0000,213.0000) -- (500.0000,211.0000) --
    (500.0000,209.0000) -- (499.0000,207.0000) -- (498.0000,206.0000) --
    (498.0000,205.0000) -- (488.0000,196.0000) -- (486.0000,193.0000);
  \path[draw,fill] (470.0000,200.0000) -- (471.0000,197.0000) --
    (475.0000,194.0000) -- (478.0000,192.0000) -- (478.0000,193.0000) --
    (475.0000,195.0000) -- (472.0000,198.0000) -- (471.0000,200.0000) --
    (470.0000,200.0000);
  \path[draw,fill] (624.0000,323.0000) -- (626.0000,322.0000) --
    (629.0000,320.0000) -- (630.0000,319.0000) -- (630.0000,318.0000) --
    (630.0000,317.0000) -- (630.0000,316.0000) -- (630.0000,315.0000) --
    (630.0000,314.0000) -- (629.0000,314.0000) -- (628.0000,313.0000) --
    (627.0000,313.0000) -- (626.0000,313.0000) -- (625.0000,313.0000) --
    (625.0000,314.0000) -- (624.0000,314.0000) -- (623.0000,314.0000) --
    (622.0000,313.0000) -- (622.0000,312.0000) -- (623.0000,311.0000) --
    (624.0000,311.0000) -- (626.0000,311.0000) -- (627.0000,311.0000) --
    (628.0000,312.0000) -- (630.0000,312.0000) -- (630.0000,313.0000) --
    (631.0000,313.0000) -- (631.0000,314.0000) -- (631.0000,315.0000) --
    (632.0000,316.0000) -- (632.0000,317.0000) -- (631.0000,318.0000) --
    (631.0000,319.0000) -- (631.0000,320.0000) -- (631.0000,321.0000) --
    (630.0000,321.0000) -- (630.0000,322.0000) -- (629.0000,322.0000) --
    (628.0000,323.0000) -- (627.0000,323.0000) -- (627.0000,324.0000) --
    (626.0000,324.0000) -- (625.0000,324.0000) -- (624.0000,323.0000);
\end{scope}
\begin{scope}[draw=black,fill=cfdae61,line join=round,line width=0.208pt]
  \path[draw,fill] (579.0000,127.0000) -- (580.0000,126.0000) --
    (583.0000,124.0000) -- (584.0000,123.0000) -- (585.0000,123.0000) --
    (587.0000,122.0000) -- (588.0000,122.0000) -- (590.0000,122.0000) --
    (592.0000,123.0000) -- (593.0000,124.0000) -- (594.0000,124.0000) --
    (596.0000,125.0000) -- (597.0000,127.0000) -- (598.0000,128.0000) --
    (599.0000,130.0000) -- (599.0000,131.0000) -- (598.0000,131.0000) --
    (598.0000,132.0000) -- (599.0000,133.0000) -- (599.0000,134.0000) --
    (598.0000,134.0000) -- (596.0000,138.0000) -- (593.0000,142.0000) --
    (590.0000,143.0000) -- (588.0000,143.0000) -- (587.0000,143.0000) --
    (586.0000,143.0000) -- (584.0000,143.0000) -- (583.0000,142.0000) --
    (581.0000,141.0000) -- (580.0000,140.0000) -- (579.0000,139.0000) --
    (577.0000,137.0000) -- (576.0000,135.0000) -- (575.0000,135.0000) --
    (575.0000,134.0000) -- (575.0000,133.0000) -- (575.0000,132.0000) --
    (575.0000,131.0000) -- (579.0000,127.0000);
  \path[draw,fill] (578.0000,136.0000) -- (589.0000,142.0000) --
    (593.0000,138.0000) -- (594.0000,137.0000) -- (595.0000,136.0000) --
    (595.0000,135.0000) -- (595.0000,134.0000) -- (596.0000,133.0000) --
    (596.0000,132.0000) -- (597.0000,131.0000) -- (597.0000,130.0000) --
    (596.0000,129.0000) -- (593.0000,127.0000) -- (590.0000,125.0000) --
    (587.0000,124.0000) -- (586.0000,124.0000) -- (586.0000,125.0000) --
    (585.0000,125.0000) -- (583.0000,127.0000) -- (580.0000,129.0000) --
    (578.0000,131.0000) -- (577.0000,132.0000) -- (577.0000,133.0000) --
    (576.0000,133.0000) -- (576.0000,134.0000) -- (577.0000,134.0000) --
    (577.0000,135.0000) -- (578.0000,135.0000) -- (578.0000,136.0000);
  \path[draw,fill] (608.0000,188.0000) -- (607.0000,188.0000) --
    (606.0000,189.0000) -- (605.0000,189.0000) -- (604.0000,190.0000) --
    (603.0000,192.0000) -- (602.0000,193.0000) -- (602.0000,194.0000) --
    (602.0000,196.0000) -- (602.0000,197.0000) -- (602.0000,198.0000) --
    (601.0000,199.0000) -- (601.0000,201.0000) -- (601.0000,202.0000) --
    (601.0000,203.0000) -- (601.0000,205.0000) -- (602.0000,206.0000) --
    (603.0000,207.0000) -- (604.0000,208.0000) -- (605.0000,209.0000) --
    (606.0000,210.0000) -- (607.0000,210.0000) -- (608.0000,210.0000) --
    (609.0000,210.0000) -- (614.0000,210.0000) -- (618.0000,211.0000) --
    (622.0000,213.0000) -- (623.0000,212.0000) -- (625.0000,212.0000) --
    (626.0000,211.0000) -- (627.0000,209.0000) -- (628.0000,208.0000) --
    (629.0000,206.0000) -- (630.0000,205.0000) -- (630.0000,204.0000) --
    (630.0000,205.0000) -- (631.0000,206.0000) -- (631.0000,207.0000) --
    (631.0000,208.0000) -- (631.0000,209.0000) -- (631.0000,210.0000) --
    (630.0000,210.0000) -- (630.0000,211.0000) -- (629.0000,212.0000) --
    (628.0000,212.0000) -- (627.0000,213.0000) -- (625.0000,214.0000) --
    (620.0000,214.0000) -- (614.0000,214.0000) -- (613.0000,214.0000) --
    (612.0000,214.0000) -- (608.0000,214.0000) -- (604.0000,213.0000) --
    (600.0000,211.0000) -- (600.0000,210.0000) -- (599.0000,210.0000) --
    (599.0000,209.0000) -- (598.0000,209.0000) -- (598.0000,208.0000) --
    (598.0000,207.0000) -- (598.0000,206.0000) -- (598.0000,205.0000) --
    (598.0000,204.0000) -- (599.0000,203.0000) -- (599.0000,202.0000) --
    (598.0000,201.0000) -- (598.0000,200.0000) -- (597.0000,199.0000) --
    (597.0000,198.0000) -- (597.0000,197.0000) -- (598.0000,196.0000) --
    (598.0000,195.0000) -- (598.0000,193.0000) -- (599.0000,191.0000) --
    (600.0000,190.0000) -- (602.0000,188.0000) -- (603.0000,187.0000) --
    (605.0000,186.0000) -- (608.0000,188.0000);
  \path[draw,fill] (534.0000,261.0000) -- (532.0000,261.0000) --
    (529.0000,263.0000) -- (526.0000,265.0000) -- (524.0000,269.0000) --
    (523.0000,270.0000) -- (521.0000,272.0000) -- (520.0000,275.0000) --
    (520.0000,276.0000) -- (524.0000,281.0000) -- (529.0000,286.0000) --
    (533.0000,290.0000) -- (538.0000,292.0000) -- (541.0000,293.0000) --
    (542.0000,293.0000) -- (544.0000,293.0000) -- (545.0000,293.0000) --
    (546.0000,293.0000) -- (547.0000,292.0000) -- (547.0000,291.0000) --
    (548.0000,290.0000) -- (548.0000,289.0000) -- (547.0000,289.0000) --
    (546.0000,288.0000) -- (546.0000,287.0000) -- (545.0000,286.0000) --
    (546.0000,285.0000) -- (547.0000,285.0000) -- (548.0000,285.0000) --
    (548.0000,286.0000) -- (549.0000,287.0000) -- (548.0000,287.0000) --
    (549.0000,288.0000) -- (550.0000,288.0000) -- (551.0000,288.0000) --
    (552.0000,288.0000) -- (552.0000,287.0000) -- (553.0000,287.0000) --
    (553.0000,283.0000) -- (553.0000,282.0000) -- (554.0000,282.0000) --
    (555.0000,282.0000) -- (555.0000,283.0000) -- (556.0000,283.0000) --
    (557.0000,284.0000) -- (557.0000,285.0000) -- (557.0000,286.0000) --
    (558.0000,286.0000) -- (557.0000,287.0000) -- (553.0000,289.0000) --
    (550.0000,292.0000) -- (548.0000,294.0000) -- (546.0000,295.0000) --
    (542.0000,295.0000) -- (538.0000,295.0000) -- (536.0000,294.0000) --
    (533.0000,293.0000) -- (528.0000,289.0000) -- (523.0000,284.0000) --
    (521.0000,283.0000) -- (519.0000,281.0000) -- (517.0000,279.0000) --
    (516.0000,276.0000) -- (516.0000,275.0000) -- (517.0000,274.0000) --
    (518.0000,272.0000) -- (519.0000,271.0000) -- (520.0000,270.0000) --
    (521.0000,269.0000) -- (522.0000,269.0000) -- (522.0000,267.0000) --
    (525.0000,264.0000) -- (528.0000,262.0000) -- (531.0000,260.0000) --
    (533.0000,259.0000) -- (534.0000,261.0000);
  \path[draw,fill] (512.0000,285.0000) -- (508.0000,288.0000) --
    (505.0000,290.0000) -- (504.0000,292.0000) -- (501.0000,295.0000) --
    (500.0000,298.0000) -- (500.0000,299.0000) -- (500.0000,300.0000) --
    (500.0000,302.0000) -- (501.0000,303.0000) -- (502.0000,304.0000) --
    (503.0000,305.0000) -- (504.0000,306.0000) -- (505.0000,306.0000) --
    (507.0000,308.0000) -- (511.0000,311.0000) -- (515.0000,313.0000) --
    (516.0000,313.0000) -- (521.0000,314.0000) -- (526.0000,314.0000) --
    (527.0000,313.0000) -- (532.0000,308.0000) -- (533.0000,307.0000) --
    (533.0000,305.0000) -- (533.0000,304.0000) -- (533.0000,303.0000) --
    (534.0000,303.0000) -- (536.0000,303.0000) -- (536.0000,305.0000) --
    (535.0000,309.0000) -- (533.0000,312.0000) -- (530.0000,314.0000) --
    (529.0000,315.0000) -- (527.0000,316.0000) -- (525.0000,317.0000) --
    (523.0000,317.0000) -- (522.0000,317.0000) -- (520.0000,317.0000) --
    (518.0000,317.0000) -- (516.0000,316.0000) -- (515.0000,316.0000) --
    (508.0000,312.0000) -- (504.0000,310.0000) -- (503.0000,309.0000) --
    (500.0000,307.0000) -- (497.0000,303.0000) -- (496.0000,301.0000) --
    (496.0000,299.0000) -- (497.0000,296.0000) -- (499.0000,293.0000) --
    (500.0000,292.0000) -- (502.0000,289.0000) -- (505.0000,286.0000) --
    (508.0000,284.0000) -- (512.0000,285.0000);
  \path[draw,fill] (338.0000,398.0000) -- (338.0000,397.0000) --
    (337.0000,396.0000) -- (338.0000,395.0000) -- (338.0000,394.0000) --
    (339.0000,393.0000) -- (340.0000,394.0000) -- (341.0000,394.0000) --
    (343.0000,395.0000) -- (345.0000,395.0000) -- (347.0000,395.0000) --
    (349.0000,394.0000) -- (350.0000,394.0000) -- (351.0000,393.0000) --
    (352.0000,392.0000) -- (353.0000,391.0000) -- (353.0000,390.0000) --
    (354.0000,388.0000) -- (354.0000,387.0000) -- (354.0000,386.0000) --
    (353.0000,384.0000) -- (353.0000,383.0000) -- (348.0000,382.0000) --
    (343.0000,381.0000) -- (342.0000,381.0000) -- (338.0000,381.0000) --
    (338.0000,380.0000) -- (337.0000,380.0000) -- (337.0000,379.0000) --
    (338.0000,379.0000) -- (338.0000,378.0000) -- (338.0000,377.0000) --
    (339.0000,377.0000) -- (339.0000,376.0000) -- (340.0000,376.0000) --
    (341.0000,376.0000) -- (342.0000,376.0000) -- (343.0000,377.0000) --
    (344.0000,377.0000) -- (346.0000,376.0000) -- (349.0000,377.0000) --
    (351.0000,377.0000) -- (353.0000,378.0000) -- (355.0000,380.0000) --
    (356.0000,381.0000) -- (357.0000,382.0000) -- (358.0000,384.0000) --
    (359.0000,386.0000) -- (359.0000,388.0000) -- (359.0000,390.0000) --
    (358.0000,392.0000) -- (357.0000,394.0000) -- (357.0000,395.0000) --
    (357.0000,396.0000) -- (357.0000,397.0000) -- (356.0000,398.0000) --
    (356.0000,399.0000) -- (355.0000,400.0000) -- (354.0000,401.0000) --
    (353.0000,401.0000) -- (352.0000,401.0000) -- (351.0000,401.0000) --
    (350.0000,401.0000) -- (349.0000,401.0000) -- (348.0000,400.0000) --
    (338.0000,398.0000);
  \path[draw,fill] (348.0000,284.0000) -- (347.0000,285.0000) --
    (347.0000,286.0000) -- (346.0000,287.0000) -- (345.0000,288.0000) --
    (345.0000,289.0000) -- (345.0000,290.0000) -- (345.0000,292.0000) --
    (345.0000,293.0000) -- (346.0000,294.0000) -- (347.0000,295.0000) --
    (348.0000,296.0000) -- (350.0000,297.0000) -- (351.0000,297.0000) --
    (353.0000,297.0000) -- (354.0000,297.0000) -- (356.0000,297.0000) --
    (357.0000,296.0000) -- (358.0000,296.0000) -- (359.0000,296.0000) --
    (360.0000,296.0000) -- (360.0000,295.0000) -- (361.0000,295.0000) --
    (361.0000,294.0000) -- (362.0000,293.0000) -- (362.0000,292.0000) --
    (362.0000,291.0000) -- (362.0000,290.0000) -- (361.0000,289.0000) --
    (361.0000,288.0000) -- (361.0000,287.0000) -- (361.0000,286.0000) --
    (360.0000,285.0000) -- (360.0000,284.0000) -- (359.0000,283.0000) --
    (358.0000,282.0000) -- (357.0000,281.0000) -- (355.0000,281.0000) --
    (354.0000,281.0000) -- (354.0000,280.0000) -- (354.0000,279.0000) --
    (355.0000,279.0000) -- (355.0000,278.0000) -- (355.0000,279.0000) --
    (356.0000,279.0000) -- (356.0000,280.0000) -- (357.0000,280.0000) --
    (359.0000,281.0000) -- (360.0000,282.0000) -- (361.0000,283.0000) --
    (362.0000,284.0000) -- (363.0000,286.0000) -- (363.0000,288.0000) --
    (363.0000,289.0000) -- (363.0000,291.0000) -- (363.0000,292.0000) --
    (363.0000,293.0000) -- (363.0000,294.0000) -- (362.0000,296.0000) --
    (361.0000,297.0000) -- (361.0000,298.0000) -- (360.0000,298.0000) --
    (358.0000,299.0000) -- (357.0000,299.0000) -- (356.0000,299.0000) --
    (354.0000,300.0000) -- (352.0000,300.0000) -- (351.0000,300.0000) --
    (349.0000,299.0000) -- (347.0000,298.0000) -- (346.0000,298.0000) --
    (345.0000,297.0000) -- (343.0000,295.0000) -- (342.0000,294.0000) --
    (341.0000,292.0000) -- (341.0000,290.0000) -- (341.0000,288.0000) --
    (341.0000,287.0000) -- (341.0000,285.0000) -- (342.0000,284.0000) --
    (343.0000,283.0000) -- (344.0000,282.0000) -- (348.0000,284.0000);
  \path[draw,fill] (348.0000,307.0000) -- (348.0000,311.0000) --
    (348.0000,318.0000) -- (348.0000,319.0000) -- (348.0000,320.0000) --
    (348.0000,321.0000) -- (349.0000,321.0000) -- (349.0000,322.0000) --
    (350.0000,322.0000) -- (351.0000,323.0000) -- (352.0000,323.0000) --
    (353.0000,323.0000) -- (354.0000,323.0000) -- (355.0000,323.0000) --
    (356.0000,323.0000) -- (357.0000,323.0000) -- (358.0000,322.0000) --
    (359.0000,321.0000) -- (359.0000,320.0000) -- (359.0000,319.0000) --
    (359.0000,318.0000) -- (360.0000,318.0000) -- (360.0000,317.0000) --
    (361.0000,317.0000) -- (361.0000,316.0000) -- (362.0000,315.0000) --
    (362.0000,314.0000) -- (362.0000,313.0000) -- (362.0000,312.0000) --
    (361.0000,311.0000) -- (360.0000,310.0000) -- (359.0000,309.0000) --
    (357.0000,309.0000) -- (356.0000,308.0000) -- (356.0000,307.0000) --
    (357.0000,306.0000) -- (358.0000,305.0000) -- (359.0000,305.0000) --
    (359.0000,306.0000) -- (361.0000,307.0000) -- (363.0000,310.0000) --
    (365.0000,313.0000) -- (365.0000,314.0000) -- (365.0000,317.0000) --
    (363.0000,321.0000) -- (362.0000,323.0000) -- (361.0000,324.0000) --
    (359.0000,326.0000) -- (358.0000,327.0000) -- (356.0000,327.0000) --
    (354.0000,327.0000) -- (351.0000,326.0000) -- (349.0000,325.0000) --
    (347.0000,324.0000) -- (346.0000,323.0000) -- (346.0000,322.0000) --
    (345.0000,317.0000) -- (344.0000,313.0000) -- (344.0000,312.0000) --
    (344.0000,310.0000) -- (345.0000,309.0000) -- (346.0000,307.0000) --
    (346.0000,306.0000) -- (348.0000,307.0000);
  \path[draw,fill] (456.0000,248.0000) -- (462.0000,248.0000) --
    (467.0000,248.0000) -- (471.0000,246.0000) -- (473.0000,248.0000) --
    (467.0000,250.0000) -- (463.0000,252.0000) -- (459.0000,251.0000) --
    (454.0000,250.0000) -- (453.0000,250.0000) -- (451.0000,248.0000) --
    (448.0000,245.0000) -- (447.0000,242.0000) -- (446.0000,240.0000) --
    (446.0000,238.0000) -- (446.0000,236.0000) -- (446.0000,233.0000) --
    (446.0000,231.0000) -- (447.0000,229.0000) -- (449.0000,227.0000) --
    (451.0000,225.0000) -- (453.0000,224.0000) -- (453.0000,223.0000) --
    (455.0000,223.0000) -- (457.0000,222.0000) -- (460.0000,222.0000) --
    (462.0000,222.0000) -- (465.0000,223.0000) -- (467.0000,224.0000) --
    (469.0000,225.0000) -- (476.0000,228.0000) -- (475.0000,234.0000) --
    (474.0000,231.0000) -- (472.0000,229.0000) -- (469.0000,228.0000) --
    (467.0000,227.0000) -- (464.0000,226.0000) -- (462.0000,226.0000) --
    (459.0000,226.0000) -- (458.0000,226.0000) -- (456.0000,227.0000) --
    (455.0000,227.0000) -- (453.0000,228.0000) -- (452.0000,230.0000) --
    (451.0000,231.0000) -- (450.0000,233.0000) -- (450.0000,235.0000) --
    (449.0000,237.0000) -- (450.0000,238.0000) -- (450.0000,240.0000) --
    (451.0000,243.0000) -- (453.0000,246.0000) -- (456.0000,248.0000);
  \path[draw,fill] (295.0000,331.0000) -- (293.0000,333.0000) --
    (294.0000,332.0000) -- (295.0000,331.0000);
  \path[draw,fill] (293.0000,333.0000) -- (293.0000,334.0000) --
    (293.0000,335.0000) -- (293.0000,337.0000) -- (293.0000,338.0000) --
    (293.0000,339.0000) -- (294.0000,340.0000) -- (294.0000,341.0000) --
    (295.0000,342.0000) -- (296.0000,343.0000) -- (297.0000,344.0000) --
    (298.0000,345.0000) -- (299.0000,345.0000) -- (301.0000,346.0000) --
    (302.0000,346.0000) -- (304.0000,346.0000) -- (305.0000,345.0000) --
    (306.0000,345.0000) -- (307.0000,344.0000) -- (308.0000,344.0000) --
    (309.0000,343.0000) -- (310.0000,342.0000) -- (311.0000,341.0000) --
    (311.0000,340.0000) -- (312.0000,338.0000) -- (312.0000,337.0000) --
    (312.0000,336.0000) -- (311.0000,335.0000) -- (311.0000,334.0000) --
    (308.0000,332.0000) -- (304.0000,331.0000) -- (300.0000,330.0000) --
    (301.0000,330.0000) -- (303.0000,329.0000) -- (305.0000,328.0000) --
    (307.0000,328.0000) -- (309.0000,328.0000) -- (311.0000,329.0000) --
    (313.0000,330.0000) -- (314.0000,331.0000) -- (315.0000,334.0000) --
    (316.0000,336.0000) -- (316.0000,338.0000) -- (316.0000,340.0000) --
    (315.0000,342.0000) -- (314.0000,344.0000) -- (313.0000,346.0000) --
    (312.0000,347.0000) -- (311.0000,348.0000) -- (309.0000,349.0000) --
    (308.0000,350.0000) -- (306.0000,350.0000) -- (304.0000,351.0000) --
    (302.0000,350.0000) -- (300.0000,350.0000) -- (299.0000,349.0000) --
    (297.0000,348.0000) -- (296.0000,347.0000) -- (295.0000,345.0000) --
    (294.0000,345.0000) -- (293.0000,345.0000) -- (292.0000,344.0000) --
    (291.0000,343.0000) -- (290.0000,342.0000) -- (290.0000,341.0000) --
    (290.0000,340.0000) -- (290.0000,338.0000) -- (293.0000,333.0000);
  \path[draw,fill] (493.0000,193.0000) -- (492.0000,193.0000) --
    (493.0000,190.0000) -- (495.0000,187.0000) -- (497.0000,184.0000) --
    (498.0000,184.0000) -- (499.0000,183.0000) -- (500.0000,183.0000) --
    (501.0000,183.0000) -- (502.0000,183.0000) -- (503.0000,183.0000) --
    (504.0000,184.0000) -- (505.0000,184.0000) -- (506.0000,185.0000) --
    (507.0000,185.0000) -- (511.0000,188.0000) -- (515.0000,191.0000) --
    (519.0000,195.0000) -- (519.0000,196.0000) -- (520.0000,197.0000) --
    (520.0000,198.0000) -- (520.0000,199.0000) -- (520.0000,200.0000) --
    (520.0000,201.0000) -- (519.0000,202.0000) -- (518.0000,203.0000) --
    (517.0000,204.0000) -- (515.0000,205.0000) -- (510.0000,209.0000) --
    (508.0000,210.0000) -- (507.0000,210.0000) -- (507.0000,209.0000) --
    (509.0000,208.0000) -- (514.0000,204.0000) -- (518.0000,200.0000) --
    (518.0000,199.0000) -- (519.0000,198.0000) -- (518.0000,197.0000) --
    (518.0000,196.0000) -- (517.0000,195.0000) -- (513.0000,191.0000) --
    (506.0000,187.0000) -- (503.0000,185.0000) -- (502.0000,185.0000) --
    (501.0000,185.0000) -- (500.0000,185.0000) -- (499.0000,185.0000) --
    (498.0000,186.0000) -- (497.0000,186.0000) -- (497.0000,187.0000) --
    (496.0000,187.0000) -- (496.0000,188.0000) -- (496.0000,189.0000) --
    (493.0000,193.0000);
  \path[draw,fill] (402.0000,185.0000) -- (402.0000,182.0000) --
    (400.0000,181.0000) -- (401.0000,181.0000) -- (403.0000,179.0000) --
    (403.0000,178.0000) -- (404.0000,177.0000) -- (405.0000,177.0000) --
    (406.0000,176.0000) -- (407.0000,176.0000) -- (407.0000,175.0000) --
    (410.0000,173.0000) -- (413.0000,171.0000) -- (415.0000,170.0000) --
    (420.0000,170.0000) -- (423.0000,170.0000) -- (426.0000,171.0000) --
    (429.0000,174.0000) -- (430.0000,175.0000) -- (432.0000,177.0000) --
    (434.0000,179.0000) -- (434.0000,181.0000) -- (434.0000,182.0000) --
    (433.0000,182.0000) -- (432.0000,183.0000) -- (431.0000,183.0000) --
    (430.0000,184.0000) -- (428.0000,185.0000) -- (427.0000,186.0000) --
    (427.0000,188.0000) -- (426.0000,189.0000) -- (423.0000,192.0000) --
    (418.0000,197.0000) -- (416.0000,199.0000) -- (412.0000,196.0000) --
    (402.0000,185.0000);
  \path[draw,fill] (407.0000,178.0000) -- (406.0000,178.0000) --
    (405.0000,178.0000) -- (405.0000,179.0000) -- (404.0000,180.0000) --
    (404.0000,181.0000) -- (404.0000,182.0000) -- (403.0000,183.0000) --
    (404.0000,183.0000) -- (406.0000,187.0000) -- (409.0000,189.0000) --
    (412.0000,191.0000) -- (412.0000,192.0000) -- (413.0000,193.0000) --
    (414.0000,194.0000) -- (414.0000,195.0000) -- (415.0000,195.0000) --
    (416.0000,195.0000) -- (417.0000,195.0000) -- (422.0000,189.0000) --
    (423.0000,189.0000) -- (424.0000,188.0000) -- (425.0000,188.0000) --
    (425.0000,187.0000) -- (426.0000,186.0000) -- (426.0000,185.0000) --
    (428.0000,181.0000) -- (429.0000,180.0000) -- (429.0000,179.0000) --
    (429.0000,178.0000) -- (428.0000,177.0000) -- (428.0000,176.0000) --
    (427.0000,176.0000) -- (427.0000,175.0000) -- (426.0000,174.0000) --
    (425.0000,173.0000) -- (424.0000,172.0000) -- (422.0000,171.0000) --
    (421.0000,171.0000) -- (420.0000,171.0000) -- (418.0000,171.0000) --
    (416.0000,171.0000) -- (415.0000,172.0000) -- (414.0000,173.0000) --
    (413.0000,174.0000) -- (407.0000,178.0000);
\end{scope}
\begin{scope}[draw=black,fill=cffffbf,line join=round,line width=0.208pt]
  \path[draw,fill] (525.0000,442.0000) -- (515.0000,436.0000) --
    (512.0000,432.0000) -- (513.0000,432.0000) -- (515.0000,432.0000) --
    (516.0000,433.0000) -- (520.0000,436.0000) -- (521.0000,437.0000) --
    (522.0000,438.0000) -- (523.0000,438.0000) -- (524.0000,439.0000) --
    (525.0000,439.0000) -- (525.0000,442.0000);
  \path[draw,fill] (507.0000,423.0000) -- (506.0000,423.0000) --
    (505.0000,423.0000) -- (504.0000,423.0000) -- (503.0000,422.0000) --
    (503.0000,421.0000) -- (502.0000,421.0000) -- (502.0000,420.0000) --
    (502.0000,419.0000) -- (502.0000,418.0000) -- (501.0000,416.0000) --
    (498.0000,412.0000) -- (496.0000,407.0000) -- (495.0000,406.0000) --
    (494.0000,403.0000) -- (492.0000,399.0000) -- (492.0000,395.0000) --
    (492.0000,393.0000) -- (492.0000,391.0000) -- (491.0000,384.0000) --
    (490.0000,383.0000) -- (490.0000,382.0000) -- (490.0000,380.0000) --
    (490.0000,379.0000) -- (490.0000,378.0000) -- (491.0000,378.0000) --
    (492.0000,378.0000) -- (493.0000,379.0000) -- (493.0000,380.0000) --
    (494.0000,381.0000) -- (493.0000,381.0000) -- (494.0000,386.0000) --
    (495.0000,390.0000) -- (494.0000,395.0000) -- (497.0000,406.0000) --
    (500.0000,410.0000) -- (505.0000,417.0000) -- (507.0000,423.0000);
  \path[draw,fill] (486.0000,367.0000) -- (486.0000,363.0000) --
    (485.0000,359.0000) -- (485.0000,354.0000) -- (486.0000,349.0000) --
    (488.0000,345.0000) -- (489.0000,341.0000) -- (491.0000,337.0000) --
    (494.0000,334.0000) -- (496.0000,331.0000) -- (496.0000,330.0000) --
    (499.0000,327.0000) -- (502.0000,324.0000) -- (504.0000,323.0000) --
    (505.0000,322.0000) -- (506.0000,321.0000) -- (507.0000,320.0000) --
    (509.0000,319.0000) -- (509.0000,317.0000) -- (512.0000,314.0000) --
    (515.0000,313.0000) -- (518.0000,313.0000) -- (518.0000,315.0000) --
    (517.0000,316.0000) -- (515.0000,317.0000) -- (514.0000,319.0000) --
    (512.0000,319.0000) -- (511.0000,320.0000) -- (510.0000,320.0000) --
    (509.0000,321.0000) -- (508.0000,322.0000) -- (507.0000,322.0000) --
    (507.0000,323.0000) -- (507.0000,324.0000) -- (505.0000,327.0000) --
    (502.0000,332.0000) -- (500.0000,334.0000) -- (498.0000,335.0000) --
    (496.0000,337.0000) -- (494.0000,339.0000) -- (493.0000,341.0000) --
    (492.0000,343.0000) -- (491.0000,346.0000) -- (491.0000,348.0000) --
    (490.0000,352.0000) -- (490.0000,358.0000) -- (490.0000,360.0000) --
    (490.0000,362.0000) -- (490.0000,368.0000) -- (490.0000,370.0000) --
    (491.0000,376.0000) -- (488.0000,376.0000) -- (488.0000,375.0000) --
    (488.0000,372.0000) -- (488.0000,370.0000) -- (487.0000,368.0000) --
    (486.0000,367.0000);
  \path[draw,fill] (521.0000,310.0000) -- (520.0000,310.0000) --
    (520.0000,309.0000) -- (520.0000,308.0000) -- (521.0000,308.0000) --
    (521.0000,307.0000) -- (521.0000,308.0000) -- (522.0000,308.0000) --
    (522.0000,309.0000) -- (522.0000,310.0000) -- (521.0000,310.0000);
  \path[draw,fill] (527.0000,305.0000) -- (526.0000,305.0000) --
    (525.0000,305.0000) -- (525.0000,304.0000) -- (525.0000,303.0000) --
    (525.0000,302.0000) -- (526.0000,302.0000) -- (526.0000,301.0000) --
    (527.0000,301.0000) -- (528.0000,301.0000) -- (532.0000,299.0000) --
    (535.0000,296.0000) -- (536.0000,294.0000) -- (541.0000,295.0000) --
    (538.0000,299.0000) -- (535.0000,301.0000) -- (531.0000,304.0000) --
    (527.0000,305.0000);
  \path[draw,fill] (557.0000,287.0000) -- (557.0000,284.0000) --
    (560.0000,283.0000) -- (565.0000,282.0000) -- (568.0000,282.0000) --
    (568.0000,284.0000) -- (564.0000,287.0000) -- (562.0000,288.0000) --
    (557.0000,287.0000);
  \path[draw,fill] (573.0000,282.0000) -- (573.0000,279.0000) --
    (576.0000,279.0000) -- (576.0000,281.0000) -- (573.0000,282.0000);
  \path[draw,fill] (590.0000,277.0000) -- (589.0000,274.0000) --
    (593.0000,274.0000) -- (597.0000,272.0000) -- (602.0000,271.0000) --
    (604.0000,272.0000) -- (602.0000,273.0000) -- (596.0000,276.0000) --
    (590.0000,277.0000);
  \path[draw,fill] (616.0000,269.0000) -- (616.0000,265.0000) --
    (622.0000,268.0000) -- (632.0000,275.0000) -- (642.0000,287.0000) --
    (651.0000,298.0000) -- (650.0000,300.0000) -- (645.0000,295.0000) --
    (639.0000,288.0000) -- (631.0000,280.0000) -- (626.0000,276.0000) --
    (622.0000,272.0000) -- (616.0000,269.0000);
\end{scope}
\begin{scope}[draw=black,fill=cabdda4,line join=round,line width=0.208pt]
  \path[draw,fill] (321.0000,481.0000) -- (324.0000,478.0000) --
    (334.0000,476.0000) -- (341.0000,472.0000) -- (353.0000,470.0000) --
    (367.0000,466.0000) -- (379.0000,462.0000) -- (386.0000,455.0000) --
    (394.0000,444.0000) -- (382.0000,466.0000) -- (375.0000,469.0000) --
    (365.0000,470.0000) -- (346.0000,475.0000) -- (321.0000,481.0000);
  \path[draw,fill] (367.0000,142.0000) -- (366.0000,145.0000) --
    (370.0000,146.0000) -- (369.0000,150.0000) -- (363.0000,147.0000) --
    (356.0000,150.0000) -- (355.0000,150.0000) -- (351.0000,151.0000) --
    (347.0000,154.0000) -- (344.0000,157.0000) -- (341.0000,161.0000) --
    (340.0000,162.0000) -- (339.0000,163.0000) -- (333.0000,171.0000) --
    (328.0000,180.0000) -- (327.0000,181.0000) -- (324.0000,181.0000) --
    (319.0000,183.0000) -- (314.0000,185.0000) -- (311.0000,187.0000) --
    (304.0000,194.0000) -- (296.0000,202.0000) -- (290.0000,208.0000) --
    (284.0000,214.0000) -- (277.0000,220.0000) -- (271.0000,228.0000) --
    (270.0000,228.0000) -- (269.0000,231.0000) -- (263.0000,246.0000) --
    (259.0000,255.0000) -- (254.0000,266.0000) -- (246.0000,283.0000) --
    (246.0000,286.0000) -- (242.0000,304.0000) -- (240.0000,322.0000) --
    (238.0000,340.0000) -- (237.0000,358.0000) -- (237.0000,369.0000) --
    (236.0000,376.0000) -- (236.0000,384.0000) -- (237.0000,392.0000) --
    (238.0000,399.0000) -- (239.0000,400.0000) -- (240.0000,407.0000) --
    (242.0000,413.0000) -- (245.0000,419.0000) -- (248.0000,425.0000) --
    (250.0000,428.0000) -- (252.0000,431.0000) -- (257.0000,439.0000) --
    (263.0000,446.0000) -- (269.0000,452.0000) -- (274.0000,457.0000) --
    (277.0000,459.0000) -- (281.0000,464.0000) -- (287.0000,467.0000) --
    (292.0000,470.0000) -- (298.0000,472.0000) -- (300.0000,473.0000) --
    (298.0000,475.0000) -- (290.0000,473.0000) -- (284.0000,470.0000) --
    (280.0000,468.0000) -- (274.0000,464.0000) -- (269.0000,459.0000) --
    (264.0000,454.0000) -- (260.0000,448.0000) -- (257.0000,444.0000) --
    (251.0000,436.0000) -- (245.0000,428.0000) -- (240.0000,419.0000) --
    (240.0000,418.0000) -- (239.0000,416.0000) -- (234.0000,396.0000) --
    (232.0000,390.0000) -- (232.0000,379.0000) -- (233.0000,366.0000) --
    (234.0000,360.0000) -- (236.0000,323.0000) -- (236.0000,318.0000) --
    (239.0000,301.0000) -- (246.0000,270.0000) -- (249.0000,267.0000) --
    (251.0000,264.0000) -- (252.0000,260.0000) -- (253.0000,256.0000) --
    (256.0000,250.0000) -- (259.0000,244.0000) -- (261.0000,238.0000) --
    (261.0000,236.0000) -- (265.0000,231.0000) -- (273.0000,221.0000) --
    (281.0000,211.0000) -- (287.0000,204.0000) -- (289.0000,202.0000) --
    (299.0000,193.0000) -- (309.0000,185.0000) -- (323.0000,175.0000) --
    (326.0000,174.0000) -- (327.0000,173.0000) -- (329.0000,170.0000) --
    (330.0000,167.0000) -- (331.0000,164.0000) -- (335.0000,161.0000) --
    (338.0000,157.0000) -- (340.0000,153.0000) -- (340.0000,152.0000) --
    (346.0000,149.0000) -- (353.0000,146.0000) -- (359.0000,144.0000) --
    (367.0000,142.0000);
\end{scope}
\begin{scope}[draw=black,fill=c2b83ba,line join=round,line width=0.208pt]
  \path[draw,fill] (332.0000,463.0000) -- (331.0000,462.0000) --
    (330.0000,460.0000) -- (329.0000,458.0000) -- (329.0000,457.0000) --
    (329.0000,455.0000) -- (330.0000,454.0000) -- (331.0000,455.0000) --
    (335.0000,457.0000) -- (339.0000,457.0000) -- (343.0000,457.0000) --
    (348.0000,456.0000) -- (352.0000,454.0000) -- (353.0000,453.0000) --
    (356.0000,452.0000) -- (364.0000,448.0000) -- (371.0000,444.0000) --
    (378.0000,439.0000) -- (385.0000,433.0000) -- (391.0000,426.0000) --
    (396.0000,419.0000) -- (397.0000,418.0000) -- (417.0000,395.0000) --
    (436.0000,372.0000) -- (438.0000,369.0000) -- (439.0000,368.0000) --
    (445.0000,361.0000) -- (452.0000,354.0000) -- (459.0000,348.0000) --
    (461.0000,347.0000) -- (466.0000,339.0000) -- (470.0000,330.0000) --
    (474.0000,321.0000) -- (477.0000,312.0000) -- (479.0000,302.0000) --
    (480.0000,300.0000) -- (480.0000,293.0000) -- (481.0000,281.0000) --
    (480.0000,269.0000) -- (480.0000,267.0000) -- (481.0000,267.0000) --
    (483.0000,264.0000) -- (484.0000,260.0000) -- (485.0000,256.0000) --
    (484.0000,253.0000) -- (484.0000,252.0000) -- (485.0000,252.0000) --
    (485.0000,251.0000) -- (485.0000,250.0000) -- (484.0000,250.0000) --
    (484.0000,249.0000) -- (483.0000,248.0000) -- (482.0000,248.0000) --
    (481.0000,248.0000) -- (480.0000,248.0000) -- (473.0000,248.0000) --
    (480.0000,245.0000) -- (487.0000,245.0000) -- (489.0000,245.0000) --
    (490.0000,247.0000) -- (491.0000,248.0000) -- (491.0000,249.0000) --
    (492.0000,251.0000) -- (492.0000,252.0000) -- (492.0000,254.0000) --
    (491.0000,256.0000) -- (489.0000,260.0000) -- (488.0000,265.0000) --
    (487.0000,271.0000) -- (487.0000,276.0000) -- (486.0000,283.0000) --
    (485.0000,294.0000) -- (484.0000,298.0000) -- (484.0000,302.0000) --
    (480.0000,315.0000) -- (475.0000,327.0000) -- (470.0000,339.0000) --
    (464.0000,351.0000) -- (457.0000,361.0000) -- (432.0000,388.0000) --
    (407.0000,415.0000) -- (397.0000,426.0000) -- (397.0000,427.0000) --
    (392.0000,433.0000) -- (387.0000,439.0000) -- (381.0000,443.0000) --
    (375.0000,447.0000) -- (369.0000,450.0000) -- (332.0000,463.0000);
  \path[draw,fill] (478.0000,238.0000) -- (488.0000,233.0000) --
    (487.0000,230.0000) -- (486.0000,226.0000) -- (486.0000,223.0000) --
    (486.0000,221.0000) -- (484.0000,211.0000) -- (481.0000,200.0000) --
    (478.0000,192.0000) -- (477.0000,190.0000) -- (474.0000,187.0000) --
    (470.0000,184.0000) -- (466.0000,182.0000) -- (464.0000,181.0000) --
    (463.0000,178.0000) -- (460.0000,173.0000) -- (456.0000,169.0000) --
    (455.0000,167.0000) -- (445.0000,166.0000) -- (427.0000,166.0000) --
    (416.0000,163.0000) -- (417.0000,163.0000) -- (432.0000,163.0000) --
    (434.0000,162.0000) -- (439.0000,162.0000) -- (445.0000,162.0000) --
    (450.0000,163.0000) -- (452.0000,163.0000) -- (458.0000,165.0000) --
    (459.0000,166.0000) -- (461.0000,169.0000) -- (467.0000,175.0000) --
    (473.0000,180.0000) -- (476.0000,183.0000) -- (478.0000,186.0000) --
    (485.0000,198.0000) -- (487.0000,203.0000) -- (488.0000,207.0000) --
    (489.0000,217.0000) -- (492.0000,226.0000) -- (492.0000,227.0000) --
    (493.0000,230.0000) -- (493.0000,234.0000) -- (493.0000,238.0000) --
    (492.0000,241.0000) -- (482.0000,241.0000) -- (478.0000,238.0000);
  \path[draw,fill] (406.0000,164.0000) -- (392.0000,166.0000) --
    (391.0000,167.0000) -- (380.0000,169.0000) -- (369.0000,172.0000) --
    (361.0000,175.0000) -- (359.0000,176.0000) -- (351.0000,180.0000) --
    (343.0000,185.0000) -- (336.0000,190.0000) -- (331.0000,195.0000) --
    (330.0000,195.0000) -- (322.0000,201.0000) -- (314.0000,208.0000) --
    (307.0000,215.0000) -- (305.0000,218.0000) -- (299.0000,228.0000) --
    (292.0000,238.0000) -- (287.0000,245.0000) -- (283.0000,248.0000) --
    (279.0000,254.0000) -- (275.0000,259.0000) -- (271.0000,265.0000) --
    (269.0000,272.0000) -- (267.0000,278.0000) -- (266.0000,280.0000) --
    (263.0000,291.0000) -- (260.0000,301.0000) -- (258.0000,306.0000) --
    (257.0000,308.0000) -- (256.0000,313.0000) -- (256.0000,317.0000) --
    (256.0000,322.0000) -- (257.0000,325.0000) -- (258.0000,336.0000) --
    (259.0000,348.0000) -- (257.0000,356.0000) -- (254.0000,367.0000) --
    (252.0000,375.0000) -- (251.0000,378.0000) -- (252.0000,384.0000) --
    (253.0000,390.0000) -- (254.0000,392.0000) -- (260.0000,411.0000) --
    (263.0000,416.0000) -- (269.0000,427.0000) -- (272.0000,432.0000) --
    (273.0000,435.0000) -- (275.0000,439.0000) -- (278.0000,442.0000) --
    (281.0000,445.0000) -- (285.0000,447.0000) -- (286.0000,448.0000) --
    (290.0000,449.0000) -- (295.0000,451.0000) -- (300.0000,451.0000) --
    (301.0000,451.0000) -- (302.0000,451.0000) -- (304.0000,451.0000) --
    (302.0000,457.0000) -- (298.0000,458.0000) -- (293.0000,458.0000) --
    (289.0000,457.0000) -- (284.0000,455.0000) -- (281.0000,453.0000) --
    (278.0000,450.0000) -- (273.0000,444.0000) -- (267.0000,436.0000) --
    (262.0000,427.0000) -- (257.0000,419.0000) -- (254.0000,409.0000) --
    (251.0000,403.0000) -- (250.0000,400.0000) -- (248.0000,395.0000) --
    (247.0000,390.0000) -- (246.0000,384.0000) -- (247.0000,379.0000) --
    (248.0000,374.0000) -- (251.0000,362.0000) -- (253.0000,349.0000) --
    (253.0000,343.0000) -- (252.0000,328.0000) -- (251.0000,317.0000) --
    (250.0000,312.0000) -- (250.0000,308.0000) -- (250.0000,303.0000) --
    (251.0000,299.0000) -- (252.0000,298.0000) -- (255.0000,292.0000) --
    (258.0000,287.0000) -- (259.0000,281.0000) -- (259.0000,280.0000) --
    (259.0000,279.0000) -- (262.0000,271.0000) -- (266.0000,263.0000) --
    (270.0000,255.0000) -- (275.0000,248.0000) -- (281.0000,242.0000) --
    (285.0000,238.0000) -- (290.0000,229.0000) -- (297.0000,220.0000) --
    (304.0000,212.0000) -- (312.0000,204.0000) -- (316.0000,200.0000) --
    (318.0000,198.0000) -- (346.0000,178.0000) -- (353.0000,174.0000) --
    (359.0000,171.0000) -- (368.0000,168.0000) -- (378.0000,165.0000) --
    (384.0000,163.0000) -- (394.0000,161.0000) -- (398.0000,160.0000) --
    (406.0000,159.0000) -- (406.0000,164.0000);
\end{scope}

                    \end{tikzpicture}
                }
                \caption[A choropleth map of \emph{Anglisano} settlement using Jenks Natural Breaks]{A choropleth map of \emph{Anglisano} settlement created using Jenks Natural Breaks algorithm on perimeters automatically derived from geometries (using default value for classes number, 5). Color difference between ditches and compounds is well-rendered.}
                \label{fig:jenks-color}
            \end{figure}

            \begin{table}
                \centering
                \begin{tabular}{ccccc}
    \toprule
    ID & shapefile\_id & perimeter & class\_n & type\\
    \otoprule
    1 & 1 & 122.792532 & 3 & compound\\
    2 & 1 & 89.612125 & 2 & compound\\
    4 & 1 & 156.587759 & 3 & compound\\
    16 & 1 & 1606.964845 & 6 & ditch\\
    \bottomrule
\end{tabular}

                \caption[Sample geometry classification results from Anglisano settlements using Jenk Natural Breaks method]{Sample results from the classification of Anglisano settlement's structures by perimeter using the Jenks Natural Breaks method. The kind of structure is saved as text attribute in the \textsf{type} column. The \textsf{shapefile\_id} column binds the geometries to the respective settlement.}
                \label{tab:jnb-results}
            \end{table}

        \subsection{Derive and measure ditches and compounds areas}
            After distinguishing ditches from compounds, efforts have been focused on geometric data deriving. One of the most useful operations when analyzing settlements is to calculate the area enclosed in ditches or compounds. Unlike some cases published in literature, facing the area calculation of rectangular shaped \emph{longhouses} \cite{spatial-south-europe}, the present study has to deal with buried, irregularly shaped, rounded compounds.

            The simplest possible approach to compound's area calculation is subtracting from the whole compound area the area occupied by the compound wall, obtaining the extension of the enclosed space. This method has some drawbacks, mostly the needing to fix the limits of the doorway. A sample result using this approach can be seen in \fref{fig:comp}. Iterating over this method on all the compounds would have generated a consistent error so this method has been discarded.

            The problem has been solved applying a common analytical geometry procedure, considering the compound's internal surface as the envelope of the compound wall's points nearest to an inner point (the compound's \emph{centroid}). Compound's inner side points have been captured as the intersection between the compound's wall and a new segment binding the centroid with a point external to the compound. Rotating this external points by \SI{1}{\degree}, the segment itself is rotated and the intersection with the innermost point of the compound's wall is registered. This process is far better explained by %\fref{fig:comp-area}

            \begin{figure}[H]
                \centering
                \subfloat[Compund's wall surface]{
                    \label{fig:comp-wall}
                    \begin{tikzpicture}[x=1mm,y=1mm,scale=0.05]
                        \begin{scope}[shift={(-190.875,-10.875)}]
  \begin{scope}[draw=black,fill=c0f7296,line join=round,line width=0.208pt]
      \path[draw,fill] (192.0000,473.0000) -- (192.0000,472.0000) --
        (191.0000,461.0000) -- (192.0000,450.0000) -- (195.0000,440.0000) --
        (199.0000,430.0000) -- (202.0000,426.0000) -- (232.0000,445.0000) --
        (303.0000,507.0000) -- (304.0000,508.0000) -- (348.0000,524.0000) --
        (422.0000,537.0000) -- (479.0000,536.0000) -- (482.0000,536.0000) --
        (544.0000,536.0000) -- (604.0000,519.0000) -- (625.0000,508.0000) --
        (640.0000,506.0000) -- (658.0000,502.0000) -- (674.0000,494.0000) --
        (689.0000,483.0000) -- (702.0000,471.0000) -- (713.0000,456.0000) --
        (720.0000,439.0000) -- (721.0000,439.0000) -- (720.0000,432.0000) --
        (722.0000,412.0000) -- (727.0000,393.0000) -- (736.0000,374.0000) --
        (747.0000,358.0000) -- (750.0000,355.0000) -- (742.0000,320.0000) --
        (702.0000,194.0000) -- (698.0000,181.0000) -- (684.0000,149.0000) --
        (664.0000,121.0000) -- (645.0000,101.0000) -- (591.0000,106.0000) --
        (579.0000,110.0000) -- (564.0000,102.0000) -- (501.0000,84.0000) --
        (498.0000,83.0000) -- (490.0000,82.0000) -- (428.0000,87.0000) --
        (377.0000,106.0000) -- (370.0000,116.0000) -- (356.0000,129.0000) --
        (341.0000,140.0000) -- (323.0000,148.0000) -- (307.0000,152.0000) --
        (301.0000,173.0000) -- (289.0000,199.0000) -- (273.0000,222.0000) --
        (260.0000,236.0000) -- (261.0000,239.0000) -- (263.0000,242.0000) --
        (264.0000,244.0000) -- (264.0000,247.0000) -- (264.0000,250.0000) --
        (263.0000,253.0000) -- (261.0000,256.0000) -- (260.0000,258.0000) --
        (258.0000,261.0000) -- (255.0000,262.0000) -- (252.0000,264.0000) --
        (250.0000,264.0000) -- (247.0000,265.0000) -- (244.0000,264.0000) --
        (241.0000,264.0000) -- (238.0000,262.0000) -- (236.0000,261.0000) --
        (233.0000,258.0000) -- (232.0000,256.0000) -- (229.0000,254.0000) --
        (224.0000,249.0000) -- (219.0000,242.0000) -- (216.0000,235.0000) --
        (214.0000,228.0000) -- (213.0000,220.0000) -- (214.0000,213.0000) --
        (216.0000,205.0000) -- (219.0000,198.0000) -- (224.0000,192.0000) --
        (227.0000,189.0000) -- (228.0000,181.0000) -- (231.0000,167.0000) --
        (238.0000,154.0000) -- (246.0000,143.0000) -- (256.0000,133.0000) --
        (268.0000,124.0000) -- (281.0000,118.0000) -- (294.0000,115.0000) --
        (326.0000,98.0000) -- (386.0000,55.0000) -- (400.0000,48.0000) --
        (478.0000,21.0000) -- (507.0000,16.0000) -- (517.0000,14.0000) --
        (551.0000,11.0000) -- (586.0000,14.0000) -- (619.0000,23.0000) --
        (651.0000,37.0000) -- (653.0000,39.0000) -- (666.0000,42.0000) --
        (678.0000,48.0000) -- (690.0000,56.0000) -- (700.0000,66.0000) --
        (708.0000,78.0000) -- (709.0000,80.0000) -- (714.0000,91.0000) --
        (718.0000,104.0000) -- (719.0000,115.0000) -- (773.0000,208.0000) --
        (775.0000,227.0000) -- (782.0000,252.0000) -- (793.0000,276.0000) --
        (808.0000,297.0000) -- (811.0000,301.0000) -- (812.0000,322.0000) --
        (808.0000,371.0000) -- (795.0000,419.0000) -- (774.0000,464.0000) --
        (752.0000,496.0000) -- (704.0000,531.0000) -- (639.0000,558.0000) --
        (581.0000,567.0000) -- (577.0000,568.0000) -- (502.0000,585.0000) --
        (495.0000,581.0000) -- (482.0000,575.0000) -- (469.0000,571.0000) --
        (455.0000,570.0000) -- (441.0000,571.0000) -- (431.0000,574.0000) --
        (396.0000,565.0000) -- (340.0000,546.0000) -- (192.0000,473.0000);
  \end{scope}
\end{scope}
\begin{scope}[shift={(-190.875,-10.875)},draw=cff0000,line width=0.800pt]
\end{scope}
 
                    \end{tikzpicture}
                }
                \hspace{0.1\textwidth}
                \subfloat[Compound's internal surface]{
                    \label{fig:comp-area}
                    \begin{tikzpicture}[x=1mm,y=1mm,scale=0.05]
                        \begin{scope}[shift={(-191.875,-81.87406)}]
  \begin{scope}[draw=black,fill=c3cad6d,line join=round,line width=0.208pt]
    \path[draw,fill] (192.0000,450.0000) -- (214.0000,213.0000) --
      (213.0000,220.0000) -- (214.0000,228.0000) -- (216.0000,235.0000) --
      (219.0000,242.0000) -- (224.0000,249.0000) -- (229.0000,254.0000) --
      (232.0000,256.0000) -- (233.0000,258.0000) -- (236.0000,261.0000) --
      (238.0000,262.0000) -- (241.0000,264.0000) -- (244.0000,264.0000) --
      (247.0000,265.0000) -- (250.0000,264.0000) -- (252.0000,264.0000) --
      (255.0000,262.0000) -- (258.0000,261.0000) -- (260.0000,258.0000) --
      (261.0000,256.0000) -- (263.0000,253.0000) -- (264.0000,250.0000) --
      (264.0000,247.0000) -- (264.0000,244.0000) -- (263.0000,242.0000) --
      (261.0000,239.0000) -- (260.0000,236.0000) -- (273.0000,222.0000) --
      (289.0000,199.0000) -- (301.0000,173.0000) -- (307.0000,152.0000) --
      (323.0000,148.0000) -- (341.0000,140.0000) -- (356.0000,129.0000) --
      (370.0000,116.0000) -- (377.0000,106.0000) -- (428.0000,87.0000) --
      (490.0000,82.0000) -- (498.0000,83.0000) -- (501.0000,84.0000) --
      (564.0000,102.0000) -- (579.0000,110.0000) -- (591.0000,106.0000) --
      (645.0000,101.0000) -- (664.0000,121.0000) -- (684.0000,149.0000) --
      (698.0000,181.0000) -- (702.0000,194.0000) -- (742.0000,320.0000) --
      (750.0000,355.0000) -- (747.0000,358.0000) -- (736.0000,374.0000) --
      (727.0000,393.0000) -- (722.0000,412.0000) -- (720.0000,432.0000) --
      (721.0000,439.0000) -- (720.0000,439.0000) -- (713.0000,456.0000) --
      (702.0000,471.0000) -- (689.0000,483.0000) -- (674.0000,494.0000) --
      (658.0000,502.0000) -- (640.0000,506.0000) -- (625.0000,508.0000) --
      (604.0000,519.0000) -- (544.0000,536.0000) -- (482.0000,536.0000) --
      (479.0000,536.0000) -- (422.0000,537.0000) -- (348.0000,524.0000) --
      (304.0000,508.0000) -- (303.0000,507.0000) -- (232.0000,445.0000) --
      (202.0000,426.0000) -- (199.0000,430.0000) -- (195.0000,440.0000) --
      (192.0000,450.0000);
  \end{scope}
  \begin{scope}[draw=black,fill=c3cad6d,line join=round,line width=0.208pt]
  \end{scope}
  \begin{scope}[draw=black,fill=c3cad6d,line join=round,line width=0.208pt]
  \end{scope}
  \begin{scope}[draw=black,fill=c3cad6d,line join=round,line width=0.208pt]
  \end{scope}
  \begin{scope}[draw=black,fill=c3cad6d,line join=round,line width=0.208pt]
  \end{scope}
\end{scope}
\begin{scope}[shift={(-191.875,-81.87406)},draw=cff0000,line width=0.800pt]
\end{scope}
 
                    \end{tikzpicture}
                }
                \caption[Deriving compound area from difference between total occupied surface and wall surface]{The area derived from the difference between compound's total occupied surface and the wall surface is not suitable for analysis, since doorstep is not set adequately.}
                \label{fig:comp}
            \end{figure}

        \subsection{Automating compounds orientation discovering}
            % compounds orientation derived from external rectangle orientation

        %\subsection{Inner/outer compounds}
\cleardoublepage
    \chapter{Results: the webGIS interface and the statistical data}

    \vspace{0.06\textheight}
    \begin{chaptersum}
        Sommario in italiano
    \end{chaptersum}

    \section{Statistical results}
        \subsection{Ditches and compound classification}
        \subsection{Derived areas}
        \subsection{Orientation discovering}
        \subsection{The web interface structure}
        \subsection{Implementation specifications}
            \subsubsection{Data processing: operator-driven vs automated GIS methods}
            \subsubsection{Scientific reproducibility: the open source approach advantage}
\cleardoublepage
    \chapter{Conclusions}

    \begin{chaptersum}
        Il lavoro presentato mostra come sia possibile automatizzare operazioni ripetitive su dati non omogenei e irregolari, usando soltanto software di analisi geografica \emph{open source}. È stata automatizzata la distinzione tra \emph{ditches} e \emph{compounds}, la derivazione di ulteriori dati geografici relativi alle loro aree, e il calcolo di dati statistici di perimetri, accessi e orientazioni, oltre che la possibilità di osservare e analizzare questi dati usando una apposita interfaccia web che comprende un webGIS.

        L'utilizzo di \emph{software libero} si è rivelato determinante, poiché ha consentito la realizzazione dell'intero sistema a costo zero, e rispetta le regole basilari del metodo scientifico, rendendo trasparente ogni passo del processo che genera nuova conoscenza da dati esistenti. Analogamente, il software presentato si è rivelato adatto a gestire una quantità non esigua di dati, che contiene un grande potenziale di conoscenza, se adeguatamente trattato. I recenti sviluppi nel trattamento di dati geografici in archeologia, geofisica e GIS hanno segnato il movimento verso analisi sempre più impegnative per quantità di dati e ampiezza dei contesti (archeologia del paesaggio) e necessitano quindi di soluzioni tecnologiche accurate ed in grado di organizzare il lavoro su larga scala.

        In questo contesto, a fronte dell'aumento di complessità dei metodi adottati, non sembra esserci uno spostamento verso tecnologie GIS più performanti: questo lavoro suggerisce una possibile soluzione al problema, adottando database spaziali (piuttosto che Shapefile), linguaggi di programmazione open source integrati nei GIS (piuttosto che soluzioni prefabbricate) e sistemi di gestione dati basati su interfacce webGIS piuttosto che desktop.

        Il lavoro presentato mostra alcune criticità minori, di facile risoluzione; tra queste, occorrerebbe risolvere problemi di calcolo nel $4\%$ delle geometrie, scrivere algoritmi più performanti per grandi quantità di dati, permettere l'esportazione dei dati geografici generati (insieme alle statistiche) e consentire il caricamento di dati geografici da fonti ulteriori.
    \end{chaptersum}

    \section{The new approach to GIS data management}
        The current work demonstrates how is possible to programmatically execute repetitive tasks when working with spatial data. This gets a particularly high value since the data are not regularly shaped or distributed. In other words, most of the algorithms presented in this work try to make their way in a set of unordered data, conducting every single case to a known process and extracting new values (geometries) from existing ones. The operation of finding regularities in natural or human phenomenon is the first step toward a better comprehension of the phenomena itself and opens towards the definition of appropriate processing methods.

        With the methods proposed in the previous chapters, a certain number of tasks have been automated: the distinction between \emph{ditches} and \emph{compounds}, the calculation of area, perimeter, access and orientation for each open structure in each settlement, and the automatic generation and plotting of results using webGIS technologies and a standard compliant export system.

        All these results have been obtained strictly using only \emph{free} and \emph{open source} geographic software (FOSS4G) and this leads to a fundamental consideration: when approaching to scientific methods, the open source advantage is evident, because it represents the highest application of the scientific methods when dealing with complex processes. The open source approach places itself in the experimental part of the scientific method, after a hypothesis has been formulated and it must be proved --- or rejected --- using the results of a \emph{reproducible} workflow, which in contemporary science is offered by a transparent, well known, \emph{atomizable} queue of operations. Any other method is a threat to the validity of the scientific process. For this reason, all the code in this work has been released using the GNU GPL v3 free software licence \cite{wiki:gpl}, and portions of the code are presented for reference in app.~\ref{app:code}.

        Another remarkable point concerns how to deal with the increasing complexity of the registered data. With the advancements in archaeology, geophysics, GIS and all the connected disciplines (anthropology, remote sensing, etc.\ ), the comprehension of a phenomenon consists in the analysis of huge datasets more often than in the past, as we are progressively widening our knowledge horizon (moving from archaeology to landscape archeology), basing the new studies on previous published results. This inevitably concerns the study of new ways of managing and analyzing data, and this work moves in this direction.
        
        Generating more data raises the necessity to have better tools, and the consequent availableness of new complex tools generates more data than in the near past. This seems a natural cycle, and could be graphically represented by a step plot. It seems that in geophysics and archaeology we are at the end of a step, since the commonly used systems to manage spatial knowledge (ESRI Shapefiles, desktop systems, etc.\ ) have reached the limit ratio with the quantity of data they can deal with; this is the reason why we need to move to new systems. Most of the tools described in this work are contained in the next step of this imaginary plot: open source technologies, webGIS, spatial databases.

    \section{Further developments}
        While the presented solution is not perfect, it is a good starting point to build new solutions to new problems and enhance our knowledge of nature. Nonetheless, some improvements are easily achievable:

        \begin{description}
            \item[fix glitches in the area drawing algorithm]\hfill\\the errors shown in \fref{sec:derived-areas} are caused algorithm miscalculations when dealing with some compounds' geometries; while the $96\%$ of successfully analyzed compounds seems a good result, this value can be easily risen to $100\%$;
            \item[write more performant code over big amount of data]\hfill\\dealing with huge datasets needs specific expedients in the code; the current code can derive areas from vector data at the speed of 8 geometries per second; any improvement on this process is desirable when the number of geometries will grow in wider contexts;
            \item[export data in a common spatially-enabled format]\hfill\\currently the export function of the software only supports CSV files (\fref{sec:webgis}), since the exportation is limited to the numerical statistic values; with small efforts, the geographical data could be exported as spatial files (like geoJSON or ESRI Shapefiles);
            \item[import data from alternative sources]\hfill\\at the current status, only compressed ESRI Shapefiles can be imported; enabling various data sources will offer more flexibility and ease the work of the analyst.
        \end{description}
\cleardoublepage
    %******************************************
    % Materiale finale
    %******************************************
    \chapter*{Acknowledgments}

    \subsection*{Free software community}
        Simply put, very little of the current work could have been realized without standing on the shoulders of giants, as Isaac Newton said. The giants in this case are the wonderful people beyond some great \emph{free as in freedom} software. My thankfulness goes to all the guys who contributed to realize the computer programs and projects I have relied on: ArchLinux, Debian GNU/Linux, VIM, \LaTeX~and TikZ, GraphViz, GNUPlot, QuantumGIS, GDAL and GEOS, OpenJump, PostGIS, SpatiaLite, TileMill, OpenStreetMap, tmux, the K Desktop Environment, Python, (geo)Django, Twitter Bootstrap, the Archaeological Recording Kit, Nginx, LeafletJS, PHPMyAdmin, MySQL, InkScape, LibreOffice.
\cleardoublepage
    \printbibliography
    \printbibliography[heading=web,category=web] 
    \cleardoublepage
    \appendix
    \addtocontents{toc}{\protect\setcounter{tocdepth}{1}}
    \addcontentsline{toc}{part}{\appendixname} 
    \chapter{Code}
\label{app:code}
\lstset{
    language=Python,
    showspaces=false,
    basicstyle=\scriptsize,
    numbers=\left,
    numberstyle=\tiny,
    stepnumber=1,
    numbersep=5pt,
    commentstyle=\color{black!60}
}

\section{Jenks Natural Breaks classification}
\label{sec:code-jenks}

    \begin{lstlisting}
    # [A] check if features for this shp already exist in helper table
    if cur_shp.feature_set.count() == HelperDitchesNumber.objects \
        .filter(shapefile_id=cur_shp_id) \
            .count():
        cur_feat = HelperDitchesNumber.objects.filter(
            shapefile_id=cur_shp_id)
    else:
        # create helping features and fill in perimeter for each
        cur_shp_geom = get_geos_geometry(cur_shp)
        cur_shp_geom.set_srid(4326)
        cur_shp_geom.transform(3857)
        for feature in cur_shp_geom:
            if feature.geom_type == 'Polygon':
                feature = MultiPolygon(feature)

            # get perimeter in Shapefile's projection
            proj_feature = feature
            proj_feature.set_srid(cur_shp_geom.srid)
            proj_feature.transform(cur_shp.proj)

            new_feat = HelperDitchesNumber(
                poly=feature,
                shapefile_id=cur_shp_id,
                perimeter=proj_feature.length)
            new_feat.save()

        cur_shp.stat_ditch_comp = cur_shp.stat_ditch_comp = True
        cur_shp.save()

        cur_feat = HelperDitchesNumber.objects.filter(
            shapefile_id=cur_shp_id)

    # [B] check if this shapefile has custom class number defined
    if cur_shp.classes:
        class_num = cur_shp.classes
    else:
        class_num = settings.GEOSTAT_SETTINGS[
            'jenk_natural_breaks_classes']

    # [C] get all perimeters as list
    perimeters = cur_feat.values_list('perimeter', flat=True)
    perim_list = []
    for x in perimeters:
        perim_list.append(x)

    # [D] calculate Jenks Natural Breaks, save in shapefile and features
    jnb_classes_list = get_jenks_breaks(perim_list, int(class_num))
    Shapefile.objects.filter(id=cur_shp_id) \
        .update(jnb=jnb_classes_list)

    # [E] fill in the class for each feature of the helper layer
    for feature in cur_feat:
        class_val = classify(feature.perimeter, jnb_classes_list)
        feature.class_n = class_val
        feature.save()
    \end{lstlisting}

\section{Get round vertexes external to compound geometry}
\label{sec:code-vertexes}

    \begin{lstlisting}
        def get_round_vertex(angle, radius,
            point_x, point_y, projection=3857, rotation=0):
            # Returns a list of GEOS points representing a circle around the point
            # (point_x, point_y) of radius `radius`, separated of an `angle` value (in
            # degrees) from the next point.
            vertex_list = []
            angle_orig = angle
            while angle <= 360:
                vertex_x = point_x + radius * cos(radians(angle))
                vertex_y = point_y + radius * sin(radians(angle))
                if rotation > 0:
                    x = ((vertex_x - point_x)*cos(rotation) - (vertex_y - point_y)
                         * sin(rotation)) + point_x
                    y = ((vertex_x - point_x)*sin(rotation) + (vertex_y - point_y)
                         * cos(rotation)) + point_y
                    vertex_list.append(Point(x, y, srid=projection))
                else:
                    vertex_list.append(Point(vertex_x, vertex_y, srid=projection))
                angle += angle_orig

            # test printing
            testpoints = MultiPoint(vertex_list)
            print(testpoints.geojson)

            return vertex_list
    \end{lstlisting}

\section{Get innermost intersection of segment and compound}
\label{sec:code-intersection}

    \begin{lstlisting}
    def check_nearest_point(coords_tuple, srid, refer_pt, defined_nearest):
        # Checks if the point generated by a tuple of coordinates (x, y) of reference
        # system ID `srid` is near to a reference point `refer_pt` more of an already
        # defined value `defined_nearest`; if true, updates the `refer_pt` list with
        # the nearest GEOS Point `refer_pt[0]` and the distance `refer_pt[1]`.
        cur_point = Point(coords_tuple, srid=srid)
        if cur_point.distance(refer_pt) < defined_nearest[1]:
            defined_nearest = [cur_point, cur_point.distance(refer_pt)]
        else:
            pass
        return defined_nearest
    \end{lstlisting}

\section{Deriving areas}
\label{sec:code-area}

    \begin{lstlisting}
    # [A] get convex hull and its centroid
    feat_convex_hull = feature.poly.convex_hull
    feat_centroid = feat_convex_hull.centroid

    # [B] feature's hull farthest point from centroid
    max_point = Point(
        feat_convex_hull.extent[2],
        feat_convex_hull.extent[3],
        srid=3857)

    radius = max_point.distance(feat_centroid)

    # get vertexes in a circle (center=centroid), every n angle
    vertexes_list = get_round_vertex(1, radius,
                                     feat_centroid.x,
                                     feat_centroid.y,
                                     3857)

    # for each point in vertex list
    for point in vertexes_list:

    # create new line between point and centroid
        line = LineString(feat_centroid, point, srid=3857)
        # line intersects geometry: get point nearest to centroid
        try:
            intersection_line = line.intersection(feature.poly)
        except GEOSException:
            pass
        if intersection_line.num_coords == 0:  # no intersection
            pass
        # intersection in 1 point
        elif intersection_line.num_coords == 1:
            area_points_list.append(Point(
                intersection_line.coords[0]))
        # intersection in 2 or more points
        elif intersection_line.num_coords >= 2:
            nearest_point = [None, 10000000]
            # intersection generates a MultiLineString (> 2 pts)
            if intersection_line.geom_type == 'MultiLineString':
                for multiline_tuple in intersection_line.tuple:
                    for coords_tuple in multiline_tuple:
                        nearest_point = check_nearest_point(
                            coords_tuple, 3857,
                            feat_centroid,
                            nearest_point)
                area_points_list.append(nearest_point[0].tuple)
            # intersection generates a LineString (2 pts)
            else:
                for coords_tuple in intersection_line.tuple:
                    nearest_point = check_nearest_point(
                        coords_tuple, 3857,
                        feat_centroid,
                        nearest_point)
                area_points_list.append(nearest_point[0].tuple)

    # close polygon, get projected area and save
    area_points_list.append(area_points_list[0])
    internal_area_polygon = Polygon(area_points_list, srid=3857)

    proj_area_polygon = internal_area_polygon
    proj_area_polygon.transform(cur_shp.proj)

    if feature.type == 'compound':
        # recognize open/closed compound
        tr = settings.GEOSTAT_SETTINGS['open_compound_treshold']
        closed_limit = 360 - ((tr * 360) / 100)
        if area_points_list.__len__() > closed_limit:
            structure_open = False
        else:
            structure_open = True
    else:
        structure_open = None
    \end{lstlisting}

\section{Find compound access and orientation}
\label{sec:code-access}

    \begin{lstlisting}
    # iterate on all open compounds
    for compound in cur_shp.helpercompoundsarea_set \
            .filter(type='compound', open=True):
        # get sides and relative lengths as dictionary
        sides = get_side_dict(compound, 3857)
        # get longest side in area polygon as a LineString
        access_linestr = max(sides, key=sides.get)

        # get access lenght as projected value
        proj_access_linestr = access_linestr
        proj_access_linestr.transform(cur_shp.proj)

        # get the centroid of the access side
        feature_centroid = compound.poly.centroid

        # get compound's farthest point from centroid
        max_point = Point(compound.poly.convex_hull.extent[2],
                          compound.poly.convex_hull.extent[3],
                          srid=3857)
        radius = max_point.distance(feature_centroid)

        # draw cardinal points around the compound every 45 degree,
        # and rotate them by 12 degree to align perpendicularly to N
        cardinal_pts = get_round_vertex(
            45,
            radius,
            feature_centroid.x,
            feature_centroid.y,
            3857,
            12)

        # create "cake slices" using cardinal points
        polygon_list = []
        for i, item in enumerate(cardinal_pts):
            points = (feature_centroid.coords,
                      item.coords,
                      cardinal_pts[i - 1].coords,
                      feature_centroid.coords)
            polygon_list.append(Polygon(points, srid=3857))
        sectors = MultiPolygon(polygon_list, srid=3857)

        # get access side centroid
        access_centroid = access_linestr.centroid
        access_centroid.transform(3857)

        # find sector containing the access centroid; get direction
        for sector in sectors:
            if sector.contains(access_centroid):
                direction = sectors.index(sector)
            else:
                pass

        # save the access LineString in a separate table
        new_compound_access = HelperCompoundsAccess(
            shapefile_id=cur_shp_id,
            comp=compound,
            poly=access_linestr,
            length=proj_access_linestr.length,
            orientation=direction
        )
        new_compound_access.save()
    \end{lstlisting}

\end{document}
